\documentclass[aop,preprint]{imsart}
\RequirePackage{amsthm,amsmath,amsfonts,amssymb}
\RequirePackage[numbers]{natbib}
\usepackage{mathtools,booktabs,longtable}
\usepackage[hidelinks]{hyperref}
\usepackage{xurl}
\startlocaldefs
\newtheorem{theorem}{Theorem}[section]
\newtheorem{proposition}[theorem]{Proposition}
\newtheorem{lemma}[theorem]{Lemma}
\newtheorem{corollary}[theorem]{Corollary}
\theoremstyle{definition}
\newtheorem{remark}[theorem]{Remark}
\newcommand{\CE}{C_{\mathrm E}}
\newcommand{\E}{\mathbb E}
\newcommand{\PP}{\mathbb P}
\newcommand{\pos}[1]{(#1)_+}
\numberwithin{equation}{section}
\endlocaldefs
\begin{document}
\hypersetup{pdfsubject={Working draft for mathematical review}}
\begin{frontmatter}
\pdfsubject{Working draft for mathematical review}
\title{Completing the iid Berry--Esseen inequality at Esseen's constant}
\runtitle{The iid Berry--Esseen inequality}
\begin{aug}
\author[A]{\fnms{Linjie}~\snm{Huang}\ead[label=e1]{2025122310@jou.edu.cn}}
\address[A]{Jiangsu Ocean University\printead[presep={,\ }]{e1}}
\end{aug}
\begin{abstract}
We give a computer-assisted completion of the iid Berry--Esseen inequality
at Esseen's constant for arbitrary standardized real distributions with
finite third absolute moment and every positive sample size. The
finite-index argument uses a common coupling of conditional residuals:
all incoming binomial count layers are charged against the loss at the
current layer. One-sided second-moment comparison bounds the resulting
test-function expectations by explicit finite reference laws, while the
third-moment cost remains that of the original distribution. The noise,
inward support and moment cost are constrained by one shared resource
budget. Convexity, lower curvature bounds and finite lower envelopes of
valid prices extend finite calculations over continuous parameter domains,
including atom activations and closed boundaries. These comparisons
complete the five actual-index implications from 26 through 30.
Combined in predecessor order with the supporting low-index, bridge
and local arguments from the moment-defect framework of
Rodziewicz and Ko\l{}odziejek, they cover all sample sizes.
A Gaussian Peano remainder supplies the initial support radii without
index subdivisions. Using the existing heat-smoothing estimate, we start the uniform local induction
at 47 from two sharp predecessors. We state
the exact global inputs and distinguish their continuous proofs from
the rational computations certifying the remaining regions.
\end{abstract}
\begin{keyword}[class=MSC]
\kwdgroup[type=primary]{\kwd{60F05}}
\kwdgroup[type=secondary]{\kwd{60E15}}
\end{keyword}
\begin{keyword}
\kwd{Berry--Esseen inequality}
\kwd{normal approximation}
\kwd{moment comparison}
\kwd{computer-assisted proof}
\end{keyword}
\end{frontmatter}
\noindent\textit{Working draft, 28 October 2026.}

\section{The problem and the proof architecture}
Let \(X_1,X_2,\ldots\) be iid real random variables with
\(\E X_1=0\), \(\E X_1^2=1\), and \(\rho=\E|X_1|^3<\infty\). Put
\[
 \Delta_n(P)=\sup_{x\in\mathbb R}
 \left|\PP\!\left(\frac{X_1+\cdots+X_n}{\sqrt n}\le x\right)-\Phi(x)\right|,
 \qquad
 \CE=\frac{3+\sqrt{10}}{6\sqrt{2\pi}}.
\]
The sharp iid conjecture asks whether
\[
 \sqrt n\,\Delta_n(P)\le\CE\rho
 \quad\text{for every such }P\text{ and every }n\ge1.
\]
Esseen's Bernoulli asymptotic lower bound makes this constant necessary;
Schulz proved the complete binomial case \cite{Schulz2016}.
The mathematical issue is the uniform bound for arbitrary real laws,
rather than the identification of the candidate constant.
An eventual-in-sample-size theorem of He and Cheng
\cite{HeCheng2026} gives a uniform upper bound beyond an explicit
large threshold. It does not settle the finite range.
The moment-defect framework of Rodziewicz and Ko\l{}odziejek
\cite{RK2026} provides a different route through finite-index
localization and propagation. Its public conditional continuation
retains finite-index premises. Our comparison argument addresses
the remaining five actual-index regions; the proof dependencies
needed for the all-index conclusion are made explicit below.

The finite-index argument below retains a single original law throughout.
Its conditional variances, inward support and third absolute moment pay
the same budget. The reference laws introduced by moment comparison
bound expectations of particular test functions; they do not replace
the original law in the statement of the theorem.
The proof uses four ranges: the supporting low-index results through
$25$, the five new implications through $30$, the actual-index bridge
through $46$, and a uniform induction for all $n\ge47$.
The last step uses a support-free Gaussian replacement estimate to
discharge the hypotheses of the existing local comparison theorem.

\section{Common residuals and moment prices}\label{mc:section}

Fix $0<p<1$, write $q=1-p$ and $v=pq$, and let $B$ be
Bernoulli with parameter $p$. Conditional on $B=i$, let $U$ have
the law of a centered variable $U_i$, where
\begin{equation}\label{mc:support}
 -O_0\le U_0\le a_0,\qquad -a_1\le U_1\le O_1,
 \qquad a_0+a_1\le\Delta<1.
\end{equation}
Here $O_i>0$ and $a_i\ge0$. Put $s_i=\E U_i^2$,
$m_i=\E U_i^3$ and $s=qs_0+ps_1$. We assume that
$Y=B+U-p$ is nonpositive on $\{B=0\}$ and nonnegative on
$\{B=1\}$. Thus $X=Y/\sqrt{v+s}$ has mean zero and variance one.
Throughout this section
\[
 \rho=\E|X|^3,\qquad R_0=v(1-2v)
\]
refer to this actual law.

Let $(B_i,\widetilde U_i)$ be independent copies of $(B,U)$, and put
$S_n=\sum_{i=1}^n(B_i+\widetilde U_i)$. For $1\le k<n$,
write $N=n-k$, $\pi_j=\binom njp^jq^{n-j}$,
$F_k=\sum_{j=0}^k\pi_j$, and
\[
 H_r=\sum_{j=k+1}^{k+r}\pi_j,\qquad d=1-\Delta.
\]
On an auxiliary probability space take independent pairs
$(U_{0i},U_{1i})$, each with the prescribed marginals, and set
$W_j=\sum_{i\le j}U_{1i}+\sum_{i>j}U_{0i}$.
Conditioning on $\sum B_i=j$ gives the same residual law $W_j$.

\begin{lemma}[One coupled staircase]\label{mc:staircase}
Let $t$ be real and choose any slope satisfying
\[
 \eta\ge\max_{1\le r\le N}\frac{H_r}{rd}.
\]
Then
\begin{equation}\label{mc:stoploss}
 \PP\{S_n\le k+t\}
 \le F_k+\eta\left\{t+\E(W_k-t-\pi_k/\eta)_+\right\}.
\end{equation}
If $h>0$ satisfies
\begin{equation}\label{mc:square-gate}
 \frac{H_r}{\pi_k}\le\frac{2rd}{h}+\frac{r^2d^2}{h^2}
 \quad(1\le r\le N),
\end{equation}
then
\begin{equation}\label{mc:clipped-square}
 \PP\{S_n\le k+t\}
 \le F_{k-1}+\frac{\pi_k}{h^2}\E(t+h-W_k)_+^2.
\end{equation}
In particular, writing $V=\E W_k^2=Ns_0+ks_1$ gives
\begin{equation}\label{mc:full-square}
 \PP\{S_n\le k+t\}
 \le F_k+\frac{\pi_k}{h^2}(V+t^2+2ht).
\end{equation}
When $\pi_{k+1},\ldots,\pi_n$ decrease, the least legal stop-loss
slope is $\pi_{k+1}/d$, and it suffices in
\eqref{mc:square-gate} to check $r=1$.
\end{lemma}
\begin{proof}
The support bounds give $W_{k+r}-W_k\ge-r\Delta$ in this coupling.
After discarding losses at counts below $k$, the CDF minus $F_k$
is at most $\E g(t-W_k)$, where
\[
 g(y)=-\pi_k\mathbf1_{\{y<0\}}
       +\sum_{r=1}^N\pi_{k+r}\mathbf1_{\{y\ge rd\}}.
\]
For every $y$, including each jump point,
$g(y)\le\eta y+\eta(-y-\pi_k/\eta)_+$.
Taking expectations proves \eqref{mc:stoploss}.
Likewise $g(y)+\pi_k\le\pi_k(y+h)_+^2/h^2$:
at the $r$th incoming jump this is precisely
\eqref{mc:square-gate}, and between jumps the right side increases.
This proves \eqref{mc:clipped-square}; replacing the positive part
by the full square gives \eqref{mc:full-square}.
Finally $H_r\le r\pi_{k+1}$ in the decreasing case.
\end{proof}

For $t<0$ and $V\le\overline V$, the choice
$h=-t+\overline V/(-t)$ in \eqref{mc:full-square}, when legal,
gives $F_k-\pi_k t^2/(\overline V+t^2)$. All these inequalities
hold for closed CDFs and allow arbitrary atoms in the conditional laws.

The moment identity, cubic endpoint inequality and shared budget below
come from the framework of Rodziewicz--Ko\l{}odziejek~\cite{RK2026};
we retain their allocation explicitly for the conditional prices.
\begin{lemma}[The actual third-moment budget]\label{mc:budget}
The standardized moment satisfies
\begin{equation}\label{mc:rho-identity}
 \rho(v+s)^{3/2}=R_0+3v(s_0+s_1)-qm_0+pm_1.
\end{equation}
If $s>0$, put $\alpha=ps_1/s$. For any positive cap
$a_0+a_1\le\delta$,
\begin{equation}\label{mc:allocation-floor}
 \rho\ge g_\delta(\alpha):=
 \frac{R_0+\{3[p+\alpha(q-p)]-\delta\}s
 +\{(1-\alpha)/\sqrt q+\alpha/\sqrt p\}^2s^2/\delta}
 {(v+s)^{3/2}}.
\end{equation}
The same allocation gives the count variance
\begin{equation}\label{mc:count-variance}
 V_k=s\left\{\frac{N(1-\alpha)}q+\frac{k\alpha}p\right\}.
\end{equation}
\end{lemma}
\begin{proof}
Expand $|Y|^3=(p-U_0)^3$ and $(q+U_1)^3$ in the two groups to
obtain \eqref{mc:rho-identity}. If $Z$ is centered, $Z\le a$ and
$\E Z^2=z$, then
$\E[(a-Z)(Z+z/a)^2]\ge0$ yields
$\E Z^3\le az-z^2/a$. Apply this to $U_0$ and $-U_1$.
Zero caps force the corresponding centered residual and variance
to vanish; omit their zero contributions. Cauchy--Schwarz gives
\[
 -qm_0+pm_1\ge
 -s\{(1-\alpha)a_0+\alpha a_1\}
 +\frac{(\sqrt q\,s_0+\sqrt p\,s_1)^2}{a_0+a_1}.
\]
Bound the first term below by $-s\delta$ and the denominator
above by $\delta$, then use
$v(s_0+s_1)=s[p+\alpha(q-p)]$ to obtain
\eqref{mc:allocation-floor}. Equation~\eqref{mc:count-variance}
follows from independence.
\end{proof}

If $p<1/2$, this floor increases with $\alpha$ for fixed $\delta$.
For any positive affine cap $\delta(\alpha)$ it is convex in
$\alpha$, since an affine square divided by a positive affine
function is convex. At $s=0$ both residuals vanish almost surely,
and \eqref{mc:rho-identity} gives the Bernoulli moment directly.

The support bounds also imply $s_i\le O_i a_i$ by averaging
$(U_i+\ell_i)(r_i-U_i)\ge0$ over the corresponding support
$[-\ell_i,r_i]$. Suppose an upstream resource bound for this law gives
\[
 g(a_0+a_1)+\lambda\{(1-\alpha)\mathfrak h_-
                         +\alpha\mathfrak h_+\}\le\mathfrak R,
 \qquad \lambda=ns,\quad g>0.
\]
Its coefficients must be supplied by the applicable resource theorem.
Combining it with $a_0+a_1\ge s_0/O_0+s_1/O_1$ yields
\begin{equation}\label{mc:signed-plane}
 \lambda\left\{(1-\alpha)
       \left(\mathfrak h_-+\frac g{nqO_0}\right)
 +\alpha\left(\mathfrak h_++\frac g{npO_1}\right)\right\}
 \le\mathfrak R.
\end{equation}
This necessary plane retains the actual allocation and its common noise.

\begin{lemma}[Classical two-moment comparison]\label{mc:two-moment}
For $a>0$, $z\ge0$, let $T_{a,z}$ have mean zero, variance $z$
and atoms $a,-z/a$, with probabilities
$z/(a^2+z),a^2/(a^2+z)$; put $T_{a,0}=0$.
If $U$ is centered, $U\le a$ and $\E U^2\le z$, then
\begin{equation}\label{mc:pinelis}
 \E(U-c)_+^2\le\E(T_{a,z}-c)_+^2\qquad(c\in\mathbb R).
\end{equation}
This comparison is preserved by independent addition. The reference
expectation increases when either $a$ or $z$ increases.
If $U_1,\ldots,U_m$ are independent, centered, bounded above by
one common $b>0$, and have total variance at most $V$, then
\begin{equation}\label{mc:bentkus}
 \E\left(\sum_i U_i-c\right)_+^2
 \le \E(Z_b(V)-c)_+^2,
 \qquad Z_b(V)=\sum_{i=1}^mT_{b,V/m,i},
\end{equation}
where the reference summands are independent.
\end{lemma}
\begin{proof}
Write $l=-z/a$, allowing $\E U^2\le z$ throughout.
If $c\le l$, then $\E(U-c)_+^2\le\E U^2+c^2\le z+c^2$,
the reference value. If $c\ge a$, both
sides vanish. For $l<c<a$, the pointwise majorant
\[
 (u-c)_+^2\le\left(\frac{a-c}{a-l}\right)^2(u-l)^2
 \quad(u\le a)
\]
has a nonnegative quadratic coefficient and equality at both reference
atoms. Its expectation increases with $\E U^2$, proving
\eqref{mc:pinelis}. This is the squared-positive-part case of
Pinelis~\cite[Lemma~3.2 and Theorem~2.1]{Pinelis2006}.
Apply the same comparison to $T_{a_1,z_1}$ with upper bounds
$a_2\ge a_1$ and $z_2\ge z_1$: this proves both parameter
monotonicities without differentiation. Conditioning on the other
independent summands proves preservation by addition. Zero variance
is the zero law; a centered variable with upper endpoint zero is
also zero and may be omitted. Jensen's inequality includes these
boundaries without assigning positive variance to $T_{0,z}$.
After replacement with common upper endpoint $b$, Bentkus's
variance equalization~\cite[Lemma~4.5]{Bentkus2004}, scaled from
$b=1$, gives \eqref{mc:bentkus}.
\end{proof}

Only test-function expectations are compared in this lemma.
The original conditional distributions need not have finite support.
For $K(V,c)=\E(Z_b(V)-c)_+^2$ one also has
\begin{equation}\label{mc:variance-derivative}
 0\le K_V\le1.
\end{equation}
There is an order proof. Both $K(V,c)$ and
$\E(c-Z_b(V))_+^2$ increase with $V$: for the latter use
$-T_{b,z}\overset d=T_{z/b,z}$ and the two parameter monotonicities.
Their sum is $V+c^2$, hence
\[
 0\le K(V_2,c)-K(V_1,c)\le V_2-V_1\qquad(0\le V_1\le V_2).
\]
The finite reference formula is $C^1$ for $V>0$, including atom
activations, so \eqref{mc:variance-derivative} follows.
Zero variance follows directly by Jensen and continuity.

For any finite reference sum $Z$ dominating the squared positive
parts of $W_k$, define
\begin{equation}\label{mc:Q}
 Q_Z(A)=\inf_{h>0}\frac{\E(Z-A+h)_+^2}{4h}.
\end{equation}
The inequality $x_+\le(x+h)_+^2/(4h)$ gives
$\E(W_k-A)_+\le Q_Z(A)$.
Write
\[
 M_j=\sum_{z_i-A+h>0}w_i(z_i-A)^j.
\]
On a fixed active set the value is $(M_2/h+2M_1+M_0h)/4$.
A positive stationary width satisfies $h^2=M_2/M_0$ and has
price $(\sqrt{M_0M_2}+M_1)/2$; its adjacent activation inequalities
must also hold. The objective is convex in $h$, with continuous
derivative at activations. If $Z\le A$, its infimum is zero.
The analogous clipped-square objective and its interior stationary width are
\[
 M_2/h^2+2M_1/h+M_0,\qquad h=-M_2/M_1\quad(M_1<0),
\]
subject also to the legal staircase gate.

\begin{lemma}[Convex reference prices]\label{mc:convex-prices}
For a fixed reference law $Z$, the price $Q_Z$ is nonincreasing,
convex and $1$-Lipschitz in $A$.
\end{lemma}
\begin{proof}
The map $(A,h)\mapsto\E(Z-A+h)_+^2/h$ is a convex perspective,
so partial minimization proves convexity of $Q_Z$.
At a positive optimizer,
$Q_Z'(A)=-(M_0+M_1/h)/2\in[-1,0]$, by
$M_1^2\le M_0M_2=M_0^2h^2$ and $M_0\le1$.
Convexity extends these bounds to all subgradients.
\end{proof}

\section{Phase interpolation and complementary prices}\label{pp:section}

For fixed actual parameters, let
\begin{equation}\label{pp:actual-residual}
 R(t)=\PP\{S_n\le k+t\}
 -\Phi\left(\frac{k+t-np}{\sqrt{npq+\lambda}}\right)-c_n\rho,
 \qquad \lambda=ns,
\end{equation}
where $0<c_n\le\CE/\sqrt n$. In particular $\rho$ and
$npq+\lambda$ remain fixed when $t$ changes.
The prices of Section~\ref{mc:section} majorize this residual;
the actual CDF itself need not be differentiable.

\begin{lemma}[Semiconvex interpolation]\label{pp:chord}
Let $E$ be continuous on $[a,b]$ and suppose
$E(t)+Mt^2/2$ is convex there, with $M\ge0$.
If $E(a)\le L$ and $E(b)\le U$, then
\begin{equation}\label{pp:chord-bound}
 E(a+y)\le L+\frac{U-L}{b-a}y
                   +\frac M2y(b-a-y),\qquad 0\le y\le b-a.
\end{equation}
The same conclusion holds if $E'$ is absolutely continuous and
$E''\ge-M$ almost everywhere, or if $E''\ge-M$ in the sense of
distributions. A supremum of functions sharing this semiconvexity
constant remains semiconvex wherever finite.
\end{lemma}
\begin{proof}
Apply the convex chord inequality to $E(t)+Mt^2/2$ and subtract
the quadratic. Each of the derivative conditions implies that
this function is convex. Supremum preserves convexity.
\end{proof}

For example, keep a fixed comparison law $Z$ and a legal slope
$\eta(p)$, and put
\[
 E(t)=F_k+\eta(p)\{t+Q_Z(t+\pi_k/\eta(p))\}
 -\Phi\left(\frac{k+t-np}{\sqrt{npq+\lambda}}\right)-c_n\rho.
\]
The reference term is convex in phase. The Gaussian contribution
has second derivative
\[
 \frac{(k+t-np)\phi((k+t-np)/\sqrt A)}{A^{3/2}},
 \qquad A=npq+\lambda.
\]
Consequently, throughout $p\in[p_l,p_u]\subset(0,1/2)$,
$\lambda\ge\lambda_l$ and $t\in[a,b]$, a valid constant is
\begin{equation}\label{pp:gaussian-curvature}
 M=\frac{(np_u-k-a)_+\phi(0)}
          {\{np_l(1-p_l)+\lambda_l\}^{3/2}}.
\end{equation}
In particular $M=0$ for phases $a\ge np_u-k$.
The same argument applies to a fixed clipped-square price,
whose positive-part-square term is convex in phase.
Uniform bounds on the two phase faces therefore control the whole
phase interval. Different valid lower bounds on the same actual
$\rho$ may be used to estimate those two faces.

Semiconvexity also permits an atom to cross an activation threshold
inside a probability interval. If $x$ is $C^2$ on that interval,
then $(x_+)^2$ is $C^1$, its derivative is absolutely continuous,
and almost everywhere
\begin{equation}\label{pp:activation-curvature}
 ((x_+)^2)'=2x_+x',\qquad
 ((x_+)^2)''=2\mathbf1_{\{x>0\}}(x')^2+2x_+x''.
\end{equation}
Thus interval bounds with $x_+\in[0,(\sup x)_+]$ and the indicator
in $[0,1]$ give a valid curvature bound across zero.
There is no singular derivative at the crossing. In the rest of this
section $Z_b(V)$ has $n$ independent summands, so the number $m$
in \eqref{mc:bentkus} is $n$. For instance, if
only the lowest atom of $Z_b(V)$ can lie below $y$, the exact identity
\begin{equation}\label{pp:one-clip}
 \E(Z_b(V)-y)_+^2
 =V+y^2-\left(1+\frac{V}{nb^2}\right)^{-n}(V/b+y)_+^2
\end{equation}
may be differentiated in this way. The condition that every other
atom is at least $y$ is still required.
For $n>1$ this identity also gives a uniform lower variance derivative
when $y\le y_{\max}<b$. Put
\[
 z_* =\frac{2n-\sqrt{2n(n-1)}}{n+1},\qquad
 m_* =1-\left(\frac{n-z_*}{n-y_{\max}/b}\right)^{n-1}
                         (2z_*-z_*^2).
\]
Then $K_V(V,y)\ge m_*$ throughout any region where at most the
lowest atom is clipped. If none is clipped, $K_V=1$.
Otherwise, with $r=1+V/(nb^2)$ and
$z=(V/b+y)/(br)\in[0,1]$, differentiation gives
\[
 K_V=1-\left(\frac{n-z}{n-y/b}\right)^{n-1}(2z-z^2).
\]
The denominator is positive, and the subtracted expression increases
with $y$. Its derivative in
$z$ has the sign of $2n-4nz+(n+1)z^2$, whose unique zero in
$[0,1]$ is $z_*$. This proves the bound, including activation
boundaries by continuity.

\begin{lemma}[A finite lower envelope of phase prices]\label{pp:mixture}
Suppose $R(t)\le E_i(t)$ for $1\le i\le r$, and that on $[0,H_i]$
the function $E_i$ is semiconvex with constant $M_i$ and has
endpoint upper bounds $L_i,U_i$. Choose
$M\ge\max_i M_i$ and $0<H\le\min_i H_i$, and write
\[
 A_i=L_i,\qquad B_i=\frac{U_i-L_i}{H_i}+\frac{MH_i}{2},
 \qquad q_i(t)=A_i+B_it-\frac M2t^2.
\]
Then, throughout $[0,H]$,
\begin{equation}\label{pp:mixed-bound}
 R(t)\le q(t):=\min_i(A_i+B_it)-\frac M2t^2.
\end{equation}
The maximum of $q$ is attained among $0,H$, the vertices $B_i/M$
when $M>0$, and the intersections $(A_j-A_i)/(B_i-B_j)$ when
$B_i\ne B_j$, retaining only points in $[0,H]$.
\end{lemma}
\begin{proof}
Lemma~\ref{pp:chord} gives $E_i\le q_i$ on its own interval.
Their minimum therefore majorizes the same actual residual. Between
consecutive intersections of the affine parts, $q$ agrees with one
concave quadratic. Its maximum is at a vertex or a segment endpoint.
\end{proof}

No mixing parameter is needed. The old fixed-mixture certificates
remain valid because, for every fixed probability vector $(\theta_i)$,
\[
 \min_iq_i(t)\le\sum_i\theta_iq_i(t).
\]
Thus any certified upper bound for such a mixture also bounds the
lower envelope. With rational face bounds and curvature, every point
in the finite candidate set and every value evaluated there is rational.
An estimate using a coefficient $c_i\le c_n$ already majorizes
the residual with coefficient $c_n$, since $\rho\ge0$.
Thus the directed lower coefficients in the finite calculation need
no alignment cost when the target is $c_n=\CE/\sqrt n$.

\begin{lemma}[A short CDF bridge]\label{pp:cdf-bridge}
For the same actual law in \eqref{pp:actual-residual}, suppose
$npq+\lambda\ge A_*>0$ and $R(b)\le U$. Then, for every $t\le b$,
\begin{equation}\label{pp:bridge-bound}
 R(t)\le U+\frac{\phi(0)}{\sqrt{A_*}}(b-t).
\end{equation}
\end{lemma}
\begin{proof}
The actual closed CDF is nondecreasing. Hence
$R(t)\le R(b)+\Phi((k+b-np)/\sqrt{npq+\lambda})
-\Phi((k+t-np)/\sqrt{npq+\lambda})$.
Bound the last difference by the maximum normal density.
\end{proof}

Only the price at $b$ is needed in this bridge. Its comparison
law, activation conditions or curvature estimates need not extend
to the smaller phases. The fixed actual moment and variance, and
monotonicity of the actual CDF, suffice.

\section{Continuous envelopes for the finite certificates}\label{ce:section}

The finite calculations use two envelopes. The first freezes a
moment-comparison law on an entire parameter box. The second keeps the
same noise in the comparison variance and moment cost until a derivative
test licenses its elimination. The tests below are obligations of each
certificate, rather than consequences asserted for every parameter choice.

\begin{proposition}[A fixed-reference box]\label{ce:fixed-box}
Let $p\in[p_l,p_u]\subset(0,1/2)$ and $\ell\le\lambda\le u$, with
$\nu=np_l(1-p_l)+\ell>0$. Suppose a reference expression $A(p,t)$,
independent of the actual noise, and a differentiable moment floor
$g(p,\lambda)$ give
\[
 R(t)\le A(p,t)-\Phi\left(\frac{k+t-np}{\sqrt{npq+\lambda}}\right)
                      -c_ng(p,\lambda).
\]
If $g_\lambda\ge b$ throughout the box, then
\begin{align}\label{ce:fixed-box-bound}
 R(t)\le{}&A(p,t)-\Phi\left(\frac{k+t-np}{\sqrt{npq+\ell}}\right)
                      -c_ng(p,\ell)+J(t),\\
 J(t)={}&(u-\ell)
 \left\{\frac{(k+t-np_l)_+\phi(0)}{2\nu^{3/2}}-c_nb\right\}_+.
 \notag
\end{align}
If $A(p,\cdot)$ is convex, the upper function has the phase
semiconvexity constant in \eqref{pp:gaussian-curvature} with
$\lambda_l=\ell$.
\end{proposition}
\begin{proof}
For $\mu=k+t-np$, differentiation of the explicit surrogate gives
\[
 \frac\partial{\partial\lambda}
 \left\{-\Phi\left(\frac\mu{\sqrt{npq+\lambda}}\right)-c_ng(p,\lambda)\right\}
 =\frac{\mu\phi(\mu/\sqrt{npq+\lambda})}{2(npq+\lambda)^{3/2}}
                           -c_ng_\lambda.
\]
Its upper bound is the expression in braces defining $J$.
Integrate from $\ell$ to the actual $\lambda$. The nonnegative
correction $J$ is convex in $t$, so only the Gaussian term contributes
to the lower phase-curvature bound.
\end{proof}

Here $A$ may be the current-count price
\[
 A(p,t)=F_k+\eta(p)\{t+Q_Z(t+\pi_k/\eta(p))\},
\]
or a fixed reflected-square reference price. For example, upper bounds
$a,z_0,z_1$ on the first group's inward support and the two conditional
variances give
\[
 Z=\sum_{i=1}^{n-k}T_{a,z_0,i}+\sum_{i=1}^{k}T_{O_1,z_1,i}.
\]
Lemma~\ref{mc:two-moment} makes this one law valid throughout the box.
The necessary resource plane \eqref{mc:signed-plane} is retained when
bounding $\lambda(1-\alpha)$ and $\lambda\alpha$ for $z_0,z_1$.
An inward cap and an allocation lower bound similarly give
\[
 a_0\le\delta-\frac{\alpha\lambda}{npO_1}.
\]
These relations keep the variance and support corners within a single
proved resource budget.

The probability direction can also be bounded without a probability
mesh. At a fixed phase, put $a(p)=t+\pi_k(p)/\eta(p)$.
For the term $\eta Q_Z(a)$, a lower second-derivative bound is supplied by
\begin{equation}\label{ce:probability-curvature}
 \eta''Q_Z(a)+2\eta'Q_Z'(a)a'+\eta Q_Z'(a)a''.
\end{equation}
The omitted term is $\eta Q_Z''(a)(a')^2\ge0$, interpreted weakly
at an activation point. The remaining binomial, Gaussian and moment-floor
derivatives are explicit. Bounds on these expressions and on
$-1\le Q_Z'\le0$ give a uniform probability chord by
Lemma~\ref{pp:chord}. Face bounds obtained from different moment
estimates are combined through the same fixed-reference envelope with
the actual law's original moment and variance.

\begin{proposition}[A coherent reflected-square envelope]\label{ce:coherent}
Let $j=k+1<n$ and fix a legal square width $h$ for reference count $j$.
Suppose the centered reflected residual summands have one common upper
bound $b>0$. Set
\[
 V_j=\frac\lambda n\left\{\frac{(n-j)(1-\alpha)}q+\frac{j\alpha}p\right\},
 \qquad K(V,y)=\E(Z_b(V)-y)_+^2.
\]
Let $G(p,\lambda,\alpha)$ be a valid actual-moment floor, convex in
$\alpha\in[0,1]$ and differentiable there. A positive affine inward
cap in \eqref{mc:allocation-floor} supplies such a floor.
Define
\begin{align}\label{ce:coherent-price}
 E(p,\lambda,\alpha,t)={}&F_k+
     \frac{\pi_j}{h^2}K(V_j,1-t-h)
 -\Phi\left(\frac{k+t-np}{\sqrt{npq+\lambda}}\right)
 -c_nG(p,\lambda,\alpha).
\end{align}
Suppose, on the whole stated probability, noise and phase domain,
\begin{equation}\label{ce:alpha-gate}
 c_nG_\alpha(p,\lambda,0)
 \ge\frac{\pi_j\lambda(j-np)_+}{npq\,h^2},
\end{equation}
and
\begin{align}\label{ce:noise-gate}
 \frac{\pi_j(n-j)}{nq h^2}K_V(V_j(p,\lambda,0),1-t-h)
 +\frac{(k+t-np)\phi((k+t-np)/\sqrt{npq+\lambda})}
                  {2(npq+\lambda)^{3/2}}\notag\\
  {}-c_nG_\lambda(p,\lambda,0)\ge0.
\end{align}
Then for every actual $\lambda\in[\ell,u]$ and allocation,
\begin{equation}\label{ce:coherent-bound}
 R(t)\le E(p,\lambda,\alpha,t)\le E(p,u,0,t).
\end{equation}
The zero-noise boundary is included by continuity.
\end{proposition}
\begin{proof}
Apply \eqref{mc:clipped-square} at count $j$ and phase $t-1$,
then apply \eqref{mc:bentkus} to the reflected residual sum.
This proves the first upper bound with the same actual third moment,
replaced only by its valid lower floor $G$.

By \eqref{mc:variance-derivative}, $0\le K_V\le1$, and
$\partial_\alpha V_j=\lambda(j-np)/(npq)$.
Convexity gives $G_\alpha(p,\lambda,\alpha)\ge G_\alpha(p,\lambda,0)$.
Thus \eqref{ce:alpha-gate} implies $E_\alpha\le0$.
On the formal zero-allocation face, \eqref{ce:noise-gate} is exactly
$\partial_\lambda E(p,\lambda,0,t)\ge0$. Integrating the two
monotonicities proves \eqref{ce:coherent-bound}.
\end{proof}

Zero allocation here is an endpoint of the majorizing function; it need
not describe an admissible original law. The first elimination uses the
universal bound $0\le K_V\le1$. Only afterwards is a one-clip formula
needed to evaluate the formal face or obtain the sharper lower derivative
in \eqref{ce:noise-gate}. Its next-atom condition is checked on that face.
The resulting explicit function $E(p,u,0,t)$ is bounded in probability
and phase by the weak-curvature rules of Section~\ref{pp:section}.

\begin{corollary}[An allocation interval]\label{ce:allocation-interval}
In the same reflected-square setting, retain one allocation interval
$[\alpha_-,\alpha_+]$. Suppose the moment floor $G$ is nondecreasing
in allocation throughout the domain. For either endpoint
$\xi\in\{\alpha_-,\alpha_+\}$ define
\begin{align}\label{ce:two-corners}
 E_{\xi}(p,\lambda,t):={}&F_k+
 \frac{\pi_j}{h^2}K(V_j(p,\lambda,\xi),1-t-h)\notag\\
 &-\Phi\left(\frac{k+t-np}{\sqrt{npq+\lambda}}\right)
 -c_nG(p,\lambda,\alpha_-).
\end{align}
Then, throughout the original probability interval, including a crossing
of $j=np$,
\begin{equation}\label{ce:variance-endpoint-max}
 R(t)\le\max\{E_{\alpha_-}(p,\lambda,t),
              E_{\alpha_+}(p,\lambda,t)\}.
\end{equation}
If $j\ge np$ everywhere, only $\xi=\alpha_+$ is needed;
if $j\le np$ everywhere, only $\xi=\alpha_-$ is needed.
If $\partial_\lambda E_\xi\ge0$ on $[\ell,u]$ for every retained
branch, replace $\lambda$ by $u$ in the resulting maximum.
\end{corollary}
\begin{proof}
The variance $V_j$ is affine in allocation, so its maximum on an
interval occurs at an endpoint. Since $K$ increases with variance
and the moment cost is at least $G(p,\lambda,\alpha_-)$,
\eqref{ce:variance-endpoint-max} follows. The derivative
$\partial_\alpha V_j=\lambda(j-np)/(npq)$ selects a single endpoint
when its sign is fixed. Each retained explicit branch keeps the
same actual $\lambda$ and can be increased to $u$ under its own
noise derivative condition. Zero noise is included by continuity.
\end{proof}

The noise gate in this corollary is explicit. With
$\nu_\xi=\{(n-j)(1-\xi)/q+j\xi/p\}/n$, it is
\[
 \frac{\pi_j\nu_\xi}{h^2}K_V(\nu_\xi\lambda,1-t-h)
 +\frac{(k+t-np)\phi((k+t-np)/\sqrt{npq+\lambda})}
               {2(npq+\lambda)^{3/2}}
 -c_nG_\lambda(p,\lambda,\alpha_-)\ge0.
\]
Probability and phase bounds are also applied to each explicit branch
before taking their maximum. A single branch across a crossing is
only an auxiliary expression; the maximum in
\eqref{ce:variance-endpoint-max} gives the actual residual bound.
Both allocation rules use one endpoint test. The affine-cap moment
floor is convex by Section~\ref{mc:section}, so
$G_\alpha(p,\lambda,\alpha_-)\ge0$ implies $G(\alpha)\ge G(\alpha_-)$.
The stronger endpoint payment in \eqref{ce:alpha-gate}, with $0$
replaced by $\alpha_-$, instead makes the full price nonincreasing.
To evaluate that endpoint, write $\delta=A-\lambda(H+r\alpha)>0$,
$B=(1-\alpha)/\sqrt q+\alpha/\sqrt p$ and
$B_\alpha=1/\sqrt p-1/\sqrt q$. Direct differentiation gives
\begin{equation}\label{ce:signed-allocation-derivative}
 G_\alpha=\frac{s}{(v+s)^{3/2}}
 \left\{3(q-p)+\lambda r+s\left(
 \frac{2BB_\alpha}{\delta}+\frac{\lambda rB^2}{\delta^2}\right)\right\}.
\end{equation}
For $p\in[p_l,p_u]\subset(0,1/2)$, $0\le\lambda\le L$,
$\delta\ge\delta_*>0$ and $B\le B_*$, a lower bound on the
braced quantity is
\begin{equation}\label{ce:signed-allocation-gate}
 3(1-2p_u)+L\min(0,r)
       +\frac{L^2\min(0,r)B_*^2}{n\delta_*^2}.
\end{equation}
This is an optional sufficient bound for the endpoint test.
The omitted $2BB_\alpha$ term is nonnegative, so a positive bound
licenses the choice for either sign of $r$.
Zero noise again follows by continuity.

\section{Global inputs and their dependencies}\label{gi:section}

For a standardized real law $P$ with finite $\rho(P)=\E|X|^3$, write
\[
 R_n(P)=\frac{\sqrt n\,\Delta_n(P)}{\rho(P)},\qquad
 C_n=\sup_P R_n(P).
\]
The supremum is over all such laws, without a support restriction.
We use the classical iid moment envelope \cite{Shevtsova2013}
\begin{equation}\label{gi:moment-envelope}
 R_n(P)\le T(\rho):=
 \min\{.469,\ .3322(1+.429/\rho),\ .3031(1+.646/\rho)\}.
\end{equation}
All terminating decimals denote rational numbers.

We include the variational starting point and contact argument;
compare He--Cheng~\cite[Lemmas~6.1 and~6.3, equations~(12)--(15)]{HeCheng2026}.
\begin{lemma}[A contrary extremizer]\label{gi:extremizer}
If $C_n\ge\CE$, the supremum defining $C_n$ is attained by a
standardized law $P$ with $\rho(P)<2$. After reflection if necessary,
there is a finite $H$ such that
\[
 F_n(H)-\Phi(H/\sqrt n)=C_n\rho(P)/\sqrt n,
 \qquad F_n(H)=P^{*n}(( -\infty,H]).
\]
\end{lemma}
\begin{proof}
The classical moment envelope in \eqref{gi:moment-envelope} gives
$C_n\le.469<1/2$ and
$T(2)\le.3031(1+.646/2)=.4010013<\CE$.
Thus a maximizing sequence whose scores tend to $C_n\ge\CE$
eventually has third moments below $2$. Tightness and this third-moment
bound give a weakly convergent subsequence $P_j\Rightarrow P$ whose
first two moments remain $0,1$. Pass also to
$\rho(P_j)\to\rho_*\in[1,2]$, so $\rho(P)\le\rho_*$.

Reflection converts a negative discrepancy into a limit of positive
closed discrepancies. Choose corresponding thresholds $H_j$ with
an error tending to zero. These thresholds stay bounded: Chebyshev's
inequality and the normal tail bound make the discrepancy tend uniformly
to zero as $|H_j|\to\infty$. Passing to $H_j\to H$, the closed-set
inequality for weak convergence gives
\[
 F_n(H)-\Phi(H/\sqrt n)\ge C_n\rho_*/\sqrt n.
\]
Maximality and $\rho(P)\le\rho_*$ force equality and
$\rho(P)=\rho_*$. The moment envelope then gives $\rho(P)<2$.

For the contact condition, put $z=H/\sqrt n$ and $m=\E[X|X|]$.
Perturb the common marginal to
$Q_\varepsilon=(1-\varepsilon)P+\varepsilon\delta_x$.
Write $\mu_\varepsilon,\sigma_\varepsilon^2$ and
$\beta_\varepsilon$ for its mean, variance and centered third
absolute moment. The threshold $H$ stays fixed in these original
coordinates. By the definition of $C_n$,
\[
 J(\varepsilon)=Q_\varepsilon^{*n}(( -\infty,H])
 -\Phi\left(\frac{H-n\mu_\varepsilon}{\sigma_\varepsilon\sqrt n}\right)
 -\frac{C_n\beta_\varepsilon}{\sqrt n\,\sigma_\varepsilon^3}
 \le0=J(0).
\]
The right derivative at zero therefore gives
\begin{align}\label{gi:contact}
 \mathcal D(x)={}&n\{F_{n-1}(H-x)-F_n(H)\}
 +\phi(z)\left\{\sqrt n x+\frac z2(x^2-1)\right\}\notag\\
 &-\frac{C_n}{\sqrt n}
 \left\{|x|^3-\frac32\rho x^2-3mx+\frac12\rho\right\}\le0.
\end{align}
Here $F_0$ is the CDF of the point mass at zero when $n=1$.
Indeed the perturbed mean and variance have derivatives $x$ and
$x^2-1$, and the centered third absolute moment has derivative
$|x|^3-\rho-3mx$. The convolution derivative is exactly
$n\{F_{n-1}(H-x)-F_n(H)\}$, preserving the common marginal.
The integral of $\mathcal D$ under $P$ is zero.
Since $\mathcal D$ is upper semicontinuous, it is zero at every
point of the support, including atoms.

\end{proof}

The heat-smoothing estimate below is already present in
Rodziewicz--Ko\l{}odziejek~\cite[Section~15]{RK2026}.
\begin{lemma}[A coarse Gaussian entry]\label{gi:gaussian-entry}
For $n\ge2$,
\[
 C_n\le C_{n-1}\sqrt{\frac n{n-1}}
       +\frac{\sqrt n}{15(n-1)^{3/2}}
       +\frac{2\sqrt n}{25(n-1)^2}.
\]
In particular, $C_{n-1}\le\CE$ and $n\ge26$ imply
\begin{equation}\label{gi:small-excess}
 n(C_n-\CE)<\frac3{10},\qquad C_n<.422.
\end{equation}
\end{lemma}
\begin{proof}
Put $v=n-1$ and $F_v(t)=\Phi(t/\sqrt v)$. Conditioning on the last
summand, replacing the first $v$ summands by a Gaussian costs at most
$C_v\rho/\sqrt v$. Taylor expansion in the last standardized summand gives
\[
 \left|\E F_v(H-X)-F_v(H)-\tfrac12F_v''(H)\right|
 \le\frac{\phi(0)\rho}{6v^{3/2}}<\frac{\rho}{15v^{3/2}}.
\]
The heat identity $\partial_vF_v=F_v''/2$ gives the matching expansion
in the variance variable, with error at most
$\|\phi'''\|_\infty/(8v^2)<2/(25v^2)$. Here Gaussian Fourier
inversion gives $\|\phi'''\|_\infty\le2/\pi<16/25$.
Combine the errors, multiply by $\sqrt n/\rho$, use $\rho\ge1$,
and take the threshold and law suprema. This is the elementary
heat-smoothing estimate also used in the supporting global route.
For the last assertion put $x=n/(n-1)$. Under the sharp predecessor,
\[
 n(C_n-\CE)\le\frac{\CE x}{\sqrt x+1}
       +\left(\frac1{15}+\frac2{25\sqrt{n-1}}\right)x^{3/2}.
\]
When $n\ge26$, $x\le26/25$, $\sqrt x<51/50$ and $\CE<.41$.
Also $x/(\sqrt x+1)\le(3x+1)/8$, since
$(\sqrt x-1)^2(3\sqrt x+1)\ge0$. Consequently the right side is less than
\[
 \frac{41}{100}\frac{3(26/25)+1}{8}
 +\left(\frac1{15}+\frac2{125}\right)\frac{26}{25}\frac{51}{50}
 =\frac{747107}{2500000}<\frac3{10}.
\]
Also $\CE+3/(10\cdot26)<.422$.
In particular, the sharp bound $C_{29}\le\CE$ already implies
$C_{30}<\CE+.01$, the coarse premise of input (II)
below. The sharp index-$30$ conclusion still needs its own proof.
\end{proof}

The finite-index completion uses the following three conclusions of the
supporting proofs.

\begin{proposition}[Global inputs]\label{gi:inputs}
The low-index companion, actual-index bridge and uniform local comparison
give the following statements. Put
$\delta_n=n\pos{C_{n-1}-C_n}$.
\begin{enumerate}
\item[(I)] $C_j\le\CE$ for $1\le j\le25$;
\item[(II)] if $C_{30}<\CE+1/100$, then $C_j\le\CE$ for
every $31\le j\le46$;
\item[(III)] for every integer $n\ge47$, the conditions
\[
 0\le n(C_n-\CE)\le7/10,\qquad \delta_n=0,
 \qquad C_n<.4247,\qquad C_{n-1},C_{n-2}<.4253
\]
cannot hold simultaneously.
\end{enumerate}
This is the zero-allowance case of the original local exclusion;
the least-failure argument below supplies exactly this condition.
The exclusion (III) is uniform in $n$; its premises will be
discharged by the Gaussian entry and predecessor induction.
\end{proposition}
The precise companion statements and their continuous proof sources
are identified in Appendix~\ref{app:companions}.

% Supporting source map, retained for preparation of the companion package.
% Low indices: work/berry-esseen-n1-base-20261024/n1_base.tex;
% work/berry-esseen-n2-public-check-20261024/collect_replay_20261024.py;
% work/berry-esseen-n3-public-check-20261024/collect_replay_20261024.py;
% work/berry-esseen-n4-analytic-check-20261025/check_analytic_20261025.py;
% work/berry-esseen-n4-public-check-20261024/check_full_closure_exact_20261025.py.
% For the low-index continuous proof map, including the actual LOW,
% conserved-identity, resource and count derivations, see
% work/berry-esseen-c5-complete-20261025/fresh-low-work-20261025/source/
% work/cutoff5_20261020/audit/ANALYTIC_DEPENDENCIES_V2.txt.
% Later low-index interfaces:
% work/berry-esseen-c5-complete-20261025/five_sum_completion_20261025.tex;
% work/berry-esseen-n6-complete-20261025/six_sum_interface_20261025.tex;
% work/berry-esseen-n7-complete-20261025/seven_sum_interface_20261025.tex;
% work/berry-esseen-n8-complete-20261025/eight_sum_interface_20261025.tex;
% work/berry-esseen-n9-complete-20261025/nine_sum_interface_20261025.tex;
% work/berry-esseen-n10-complete-20261025/n10-interface/main.tex;
% work/berry-esseen-n11-14-complete-20261026/conditional_11_14_interface_20261026.tex;
% work/berry-esseen-n15-19-complete-20261026/n15-19-interface/main.tex;
% work/berry-esseen-n20-complete-20261025/twenty_sum_interface_20261025.tex;
% work/berry-esseen-n20-singlebatch-20261026/actual20_singlebatch_interface_20261026.tex;
% work/berry-esseen-n21-24-complete-20261026/n21-24-interface/main.tex;
% work/berry-esseen-n25-full-parent-20261026/n25-interface/main.tex;
% work/berry-esseen-n25-full-parent-20261026/literal25-fourier-proof/main.tex.

The classical inputs are Prawitz's smoothing inequality \cite{Prawitz1972} in the form
given by Shevtsova \cite[Lemma~4.1, p.~263]{Shevtsova2012},
Shevtsova's finite iid moment bounds \cite{Shevtsova2013}, and
Schulz's strict bound for every binomial law \cite[Theorem~1]{Schulz2016}.
In particular, the moment envelope \eqref{gi:moment-envelope} gives
$T(\rho)\le819/2000<\CE$ for $\rho\ge7361/4000$,
excluding the entire larger-moment ray.
Prawitz smoothing is applied at each actual integer, retaining its
endpoint terms and infinite Gaussian tail. The contact, moment-defect,
resource and convolution-count reductions are proved in the companions;
they are not additional unnamed literature assumptions.

Lemma~\ref{gi:extremizer} supplies the common starting point on a
contrary branch, before any use of higher moments. At this same law,
\begin{equation}\label{gi:predecessor}
 n\pos{R_{n-1}(P)-C_n}\le n\pos{C_{n-1}-C_n}.
\end{equation}
Thus a global predecessor bound licenses contact allowances throughout
the argument. A score bound restricted to a moment band is never
substituted for that predecessor.

For (I), the first four direct bounds are
\[
 C_1\le\tfrac25,\quad C_2\le\tfrac{2\sqrt2}{7},\quad
 C_3\le\tfrac{59\sqrt3}{250},\quad C_4\le.4097<\CE.
\]
The companion compositions use the unrestricted predecessor in order
through $24$. At $25$, a support-free Fourier bound excludes
$\rho\ge6/5$; the two lower bands use the global $C_{24}<.419$
and their complete resource/count covers. Their continuous reductions,
unseeded resource steps and closed arithmetic covers are the precise
inputs listed in Appendix~\ref{app:companions}. The fifth-index proof
is also given separately in \cite{RKFive2026}; subsequent source
developments are collected in \cite{RKSources2026}.

% Actual-index bridge sources:
% work/berry-esseen-n31-primary-check-20261025/sharp31_bridge_composition_20261025/main.tex;
% work/berry-esseen-n35-complete-20261025/actual35_conditional_composition_20261025.tex;
% work/berry-esseen-n36-38-complete-20261026/actual36_38_conditional_composition_20261026.tex;
% work/berry-esseen-n39-complete-20261026/actual39_two_step_composition_20261026/main.tex;
% work/berry-esseen-n40-route-20261026/n40-48-interface/main.tex.

For (II), the coarse bound at $30$ gives allowance $31/100$ at
a contrary index-$31$ extremizer. The eight closed moment bands have
actual-$31$ support, resource and count proofs. The bridges through
$34$ keep $mD_m\le31d$, acquire support at each actual $m$, and use
that integer's binomial law. Subsequent implications proceed through
$35$--$39$ and $40$--$46$, retaining both deleted predecessors where
required. Thus the bridge starts from the stated coarse premise
without any monotonicity assumption on $C_n$.

% Large-index source:
% work/berry-esseen-n49-check-20261025/n49_completion_20261025.tex;
% substantive saved-output joins:
% work/berry-esseen-n49-check-20261025/join_n49_fresh_results_20261025.py;
% work/berry-esseen-n49-check-20261025/check_n49_final_join_20261025.py.

For (III), the moment envelope gives $\rho<7361/4000$ at a contrary
extremizer. Corollary~\ref{gc:tail-support} then gives radius $3.84$
directly from contact and the deleted-predecessor allowance.
It also supplies the two left radial contractions and the central
radius-$2.5$ contraction, in their original acquisition order.
Thus the six initial index-band and subsequent radial support
calculations are replaced by the same cubic bound. The bounded-threshold
coordinate-strip exclusions retain their original proofs.
The resource and confinement schedule then supplies $45$ central
and $38$ left probability intervals, covering $[.304,.47]$ after
all complementary intervals have been excluded. Their resources
retain the actual weighted quantities $nD$ and $n\E U^2$.
The complete local comparison is uniform for every $n\ge47$.
At positive noise it gives a strict loss relative to this law's
own Bernoulli score; at zero noise Schulz's strict binomial theorem
applies directly. The acquisition includes both deleted-score caps
and the full excess interval $0\le n(C_n-\CE)\le.7$.

\begin{proposition}[Induction from two sharp predecessors]\label{gi:tail}
Input~(III) and $C_{45},C_{46}\le\CE$ imply
$C_n<\CE$ for every integer $n\ge47$.
\end{proposition}
\begin{proof}
If a failure exists, choose the least $n\ge47$ with $C_n\ge\CE$.
Both predecessors satisfy $C_{n-1},C_{n-2}\le\CE$, so $\delta_n=0$
and both deleted-score caps in (III) hold. The Gaussian entry gives
\[
 0\le n(C_n-\CE)<3/10<7/10,\qquad
 C_n<.422<.4247.
\]
All premises of (III) are now satisfied, a contradiction.
\end{proof}

The moment-defect method and the original finite subdivisions retain
their attribution to Rodziewicz--Ko\l{}odziejek \cite{RK2026}.
The conditional continuation in Section~11 of that working manuscript
retains analytic and historical certificate premises, including the
sharp assertions at $25$--$30$. It does not by itself establish (I) or
fill the intervening finite range. Inputs (I)--(III) here refer to the
companion proofs with their continuous reductions, complete exact
calculations and closed-domain compositions. In particular, the proof
order is (I), the new actual-index steps $26$--$30$, (II) through
$46$, and Proposition~\ref{gi:tail}. The five finite ceilings in
the independent $n\ge49$ route are unnecessary for this composition.

\section{From a global extremizer to a conditional comparison}\label{loc:section}

Fix an actual integer $26\le n\le30$ and assume the unrestricted
predecessor bound $C_{n-1}\le\CE$. Suppose $C=C_n\ge\CE$, and take
the positive closed maximizing pair $(P,H)$ of
Lemma~\ref{gi:extremizer}, with $1\le\rho\le6/5$.
All quantities below belong to this one pair. Write
$\beta=\E[X|X|]$, $\gamma=\E X^3$, $M_j=\E|X|^j$,
$z=H/\sqrt n$ and $\varphi=\phi(z)$.
The moment-defect construction and its original finite subdivisions are
due to Rodziewicz--Ko\l{}odziejek~\cite{RK2026}.
We spell out the continuous implications needed to interpret the finite
rows; the arithmetic verification of a row is a separate obligation.

First, the predecessor bound has a pointwise support consequence.
Set $s_n=\sqrt{n/(n-1)}$, $e_n=n(s_n-1)-1/2$, and
\[
 \mathcal H_n(H,x)=n^{3/2}
 \{\Phi((H-x)/\sqrt{n-1})-\Phi(H/\sqrt n)\}
 +nx\varphi+\tfrac H2(x^2-1)\varphi.
\]
Multiplying \eqref{gi:contact} by $\sqrt n$ and using
$F_{n-1}(H-x)\le\Phi((H-x)/\sqrt{n-1})+
C_{n-1}\rho/\sqrt{n-1}$ gives
\begin{equation}\label{loc:support-test}
 \sqrt n\,\mathcal D(x)\le\mathcal H_n(H,x)
 +C\{-|x|^3+\tfrac32\rho x^2+3\beta x+\rho e_n\}.
\end{equation}
The otherwise present allowance
$n\rho s_n(C_{n-1}-C)_+$ is zero. Since $\mathcal D=0$ on the
full support, a strictly negative upper bound for the right side on
$x\le-R_-$ or $x\ge R_+$ excludes that entire ray, including atoms.
The bound must hold for every real $H$ and every moment and score value
in the incoming intervals. Thus each support refinement uses the actual
integer $n$ and previously established moment intervals. In particular,
a support calculation at $25$ cannot simply be relabelled at $n$.
For the initial radius no such calculation is needed:
Proposition~\ref{gc:initial-support} gives $12/5$ when $\rho\le11/10$
and $13/5$ when $\rho\le6/5$, by one Peano-kernel argument for all
$n\ge25$. This precedes all higher-moment and defect refinements.

Here is the physical resource behind the refinement. Put
$\sigma=\sqrt{10}$ and $\tau(t)=\E[X\mathbf1_{\{X\ge t\}}]$.
For every bounded component $g=(a,b)$ of the complement of the closed
support, $\tau$ is constant there; call the value $\tau_g$. Define
\begin{align*}
 J&=\sum_g\tau_g(b-a)^2,&
 J_1&=\sum_g\tau_g\int_a^b(t-a)t\,dt,\\
 J_2&=\sum_g\tau_g\int_a^b(t-a)t|t|\,dt,&
 D&=(3+\sigma)\rho-\gamma-3J.
\end{align*}
All these expressions are finite by the bounded support already proved.
For $\nu=\mathcal L(H-X_1-\cdots-X_{n-1})$, also put
\[
 \mathcal L=\sum_g\tau_g\int_{[a,b)}(t-a)\,\nu(dt)\ge0.
\]
The half-open interval specifies the convention at atoms.

\begin{lemma}[The physical defect]\label{loc:defect}
One has $D\ge0$. If $D<(\sigma-1)\rho$, there is an actual central
gap $(-A_-,A_+)$ with $A_-,A_+>0$. Put
\[
 p=\PP(X>0),\quad q=1-p,\quad v=pq,\quad
 \mu=\E X_+=\E X_-,\quad
 \bar A_-=\mu/q,\quad\bar A_+=\mu/p,
\]
and $\Delta_\pm=\bar A_\pm-A_\pm\ge0$.
Let $\theta_\pm=\E[(X-\E[X\mid\operatorname{sgn}X])^2;
\ \pm X>0]$ and $\theta=\theta_-+\theta_+$.
Then, with $p_*=(4-\sigma)/2$,
\begin{align}\label{loc:joint-defect}
 D\ge{}&\frac{(6+2\sigma)(1-\theta)^{3/2}}{\sqrt v}
                    (p-p_*)^2\notag\\
 &+6\mu\{(A_-+A_+)(\Delta_-+\Delta_+)+\Delta_-\Delta_+\}
 +3(\sigma+3)A_-\theta_-+3(\sigma+1)A_+\theta_+.
\end{align}
\end{lemma}
\begin{proof}
If zero lies in the support, all gaps have one sign and
$J\le2\int|t|\tau(t)\,dt=\rho$. Hence
$D\ge\sigma\rho-\gamma\ge(\sigma-1)\rho$.
Otherwise let $J_o=J-\mu(A_-+A_+)^2$ and
\[
 T_\pm=\E X_\pm^3-\mu A_\pm^2,
 \qquad R_\pm=\E[X_\pm(X_\pm-A_\pm)].
\]
On a positive outer gap $(a,b)$,
$(b-a)^2\le2\int_a^b(t-A_+)dt$.
Summing and applying Tonelli, and then reflecting for the negative
side, proves $J_o\le T_++T_--2A_+R_+-2A_-R_-$.
Expansion gives the nonnegative decomposition
\[
 D=\mu(\sqrt{\sigma+1}A_--\sqrt{\sigma-1}A_+)^2
 +(\sigma-1)T_++(\sigma+1)T_-+3(T_++T_--J_o).
\]
For the stronger assertion, conditional centering gives
\[
 \E X_\pm^3=\mu\bar A_\pm^2+
 \E[(X_\pm+2\bar A_\pm)(X_\pm-\bar A_\pm)^2;
                         \ \pm X>0]
 \ge\mu\bar A_\pm^2+3A_\pm\theta_\pm,
\]
and $R_\pm=\theta_\pm+\mu\Delta_\pm$.
Insert these inequalities into the preceding outer-gap bound.
Use $\mu^2=(1-\theta)pq$ and
$(\sigma+1)p^2+(\sigma-1)q^2-6pq
=(6+2\sigma)(p-p_*)^2$ to obtain \eqref{loc:joint-defect}.
\end{proof}

The upper bound for this same $D$ comes from contact and Fourier
inversion, rather than from an assumed two-point approximation.
For $w=H^2/(2n)$ the exact conserved identity is
\begin{equation}\label{loc:conserved}
 n(C-\CE)\rho+n\CE\rho(1-e^{-w})
 +\frac{n\phi(0)e^{-w}}6D+n^{3/2}\mathcal L=\mathcal R,
\end{equation}
where, setting
\[
 K_1=\E[X|X|^3]-3\beta-\tfrac32\rho\gamma,\quad
 K_q=M_5/4-\rho M_4/2-\tfrac32\beta\gamma+\rho/2,
 \quad U=J_1+(M_4-3)/6,
\]
the right side is
\begin{equation}\label{loc:raw}
 \mathcal R=C(3\beta J/2+3\rho J_1-3J_2-K_q)
 +H(CK_1+\varphi U)-H^2(C\rho+\gamma\varphi/2)+E_n.
\end{equation}
The remainder has the exact representation
\begin{equation}\label{loc:fourier-remainder}
 E_n=\frac{2\sqrt n}{\pi}\int_0^\infty
 \frac{\Im\{e^{-itH}[f(t)^n-e^{-nt^2/2}(1-in\gamma t^3/6)]\}}
 {t^3}\,dt,\qquad f(t)=\E e^{itX}.
\end{equation}
This integral is absolutely convergent.
To see the identity, write $\nu([x,\infty))\le\mathfrak q(x)$ for
the cubic majorant furnished by \eqref{gi:contact}, with equality on
the support. Evaluate
$\int\tau(x)[\mathfrak q(x)-\nu([x,\infty))]dx$ first gap by gap
and then using the density $\tau$ of the zero-bias law.
The first evaluation is
$\varphi J/(2\sqrt n)+G/n^{3/2}-\mathcal L$, where
$G=3C\beta J/2+(H\varphi+3C\rho)J_1-3CJ_2$.
For the second, contact multiplied by $X$ supplies the first truncated
moment; the payoff $(x-z)|x-z|$ supplies the second truncated moment.
Their difference gives
\begin{align*}
 C\rho+\sqrt n\mathcal L
 &=\frac{\varphi}{6}(\gamma+3J)+(G-P_0)/n+E_n/n,\\
 P_0&=H^2(C\rho+\gamma\varphi/2)
      -H\{CK_1+\varphi(M_4-3)/6\}+CK_q.
\end{align*}
Substitute $\gamma+3J=(3+\sigma)\rho-D$ to obtain
\eqref{loc:conserved}--\eqref{loc:raw}.
Formula~\eqref{loc:fourier-remainder} follows from
$\int_0^\infty(ty-\sin ty)t^{-3}dt=\pi y|y|/4$ by compensated
Fubini; bounded support controls the origin and $t^{-3}$ controls
infinity. No discontinuous distribution function is inverted.

There is a necessary order in using this identity. A first bound
$\mathcal R<r<n c_0$, with $C\rho\ge c_0>0$, proves
$\varphi(\gamma+3J)/6=C\rho+\sqrt n\mathcal L-\mathcal R/n>0$.
Only then may the Gaussian weight be removed, giving
\begin{equation}\label{loc:unweighted}
 n(C-\CE)\rho+\kappa nD+n^{3/2}\mathcal L\le\mathcal R,
 \qquad \kappa=997/15000<\phi(0)/6.
\end{equation}
Every later raw bound $\mathcal R\le r'$ also gives
$e^{-w}\ge1-5r'/(2n)$ and $nD\le r'/\kappa$.
The initial pass therefore requires no defect bound of its own output.
The actual central gap is used only after
$D<(\sigma-1)\rho$ has been established.

On a closed moment cell, the finite calculation bounds
\eqref{loc:fourier-remainder} by a coherent expression
$K_F+s_F|H|$, and bounds the coefficients of \eqref{loc:raw}.
It thereby proves an envelope
$\mathcal R\le\max_j(A_j+F_{j,\pm}|H|-Q_j^{\rm raw}H^2)$,
where every $Q_j^{\rm raw}>0$.
Square completion bounds the envelope on the whole half-line.
The signed curvature uses the actual $\gamma$: if
$\gamma\ge g_0$ and $L_\phi\le\varphi\le U_\phi$, its valid price is
$C\rho+\tfrac12\min(L_\phi g_0,U_\phi g_0)$.
Reflecting an auxiliary Fourier integral does not change this sign.
Keeping the phase cost in \eqref{loc:conserved} further gives
\begin{equation}\label{loc:phase-budget}
 \kappa nD\le\max_j(A_j+F_{j,\pm}|H|-Q_jH^2).
\end{equation}
For example, a previously proved $\mathcal R\le r$ and $D\le d$
allow the additional curvature
$(1-5r/(2n))( .4095\rho-d/15)/2$, when positive.
It is included once in $Q_j$.
Thus a bound on the right side of \eqref{loc:phase-budget} is converted
to physical $D$ using $15000/(997n)$ at the actual integer.

We now give the canonical parameters required by
Section~\ref{mc:section}. Once the central gap is licensed, put
\begin{equation}\label{loc:canonical}
 B=\mathbf1_{\{X>0\}},\qquad
 \ell=\bar A_-+\bar A_+=\mu/v,\qquad
 U=X/\ell+p-B.
\end{equation}
Then $\E[U\mid B]=0$ and $X=\ell(B+U-p)$ exactly.
With $s_i=\E[U^2\mid B=i]$, $s=qs_0+ps_1$ and $\lambda_n=ns$,
\begin{equation}\label{loc:canonical-identities}
 \ell^2=\frac{1-\theta}{v}=\frac1{v+s},\qquad
 \mu\ell=1-\theta,\qquad
 \theta=\frac{s}{v+s},\qquad
 \alpha=\frac{ps_1}{s}=\frac{\theta_+}{\theta}\quad(s>0).
\end{equation}
In particular, the moment in Section~\ref{mc:section} is this law's
actual $\rho$. If its physical support is in $[-R_-,R_+]$, take
\begin{equation}\label{loc:conditional-support}
 a_0=\Delta_-/\ell,\quad a_1=\Delta_+/\ell,\quad
 O_0=R_-/\ell-p,\quad O_1=R_+/\ell-q.
\end{equation}
These give precisely $-O_0\le U_0\le a_0$ and
$-a_1\le U_1\le O_1$. The gap yields $a_0+a_1<1$.
Any zero outer cap can be enlarged to a positive cap when a comparison
formula requires division by it.
There is no probability symmetry in this construction: the original
positive discrepancy determines the orientation of $p$.

For explicit resource coefficients, suppose the already acquired cell
has $p\in I=[l,u]\subset(0,1)$, $\theta\le\Theta<1$,
$A_-\ge A_l>0$, $A_+\ge B_l>0$. Let
\begin{gather*}
 v_* =\max_{p\in I}p(1-p),\quad S_*=A_l+B_l,\\
 b=\frac{(6+2\sigma)(1-\Theta)^{3/2}}{\sqrt{v_*}}
                      \operatorname{dist}(p_*,I)^2,\quad
 g=6(1-\Theta)S_*,\\
 K_-=3(\sigma+3)A_l,\quad K_+=3(\sigma+1)B_l,\quad
 h_{n,\pm}=\frac{K_\pm(1-\Theta)}{nv_*}.
\end{gather*}
Dropping the nonnegative product in \eqref{loc:joint-defect} and
using \eqref{loc:canonical-identities} proves the single inequality
\begin{equation}\label{loc:affine-resource}
 D\ge b+g(a_0+a_1)
       +\lambda_n\{(1-\alpha)h_{n,-}+\alpha h_{n,+}\}.
\end{equation}
It supplies the upstream prices in \eqref{mc:signed-plane}, with
$\mathfrak R=\overline D-b$ whenever
$D\le\overline D$. Replacing $K_\pm$ by their minimum gives the
scalar noise price used by the finite rows. Separate upper bounds for
inward displacement and noise do not imply \eqref{loc:affine-resource}.

The stronger signed prices used in the new comparison are finite
consequences of this inequality. Write $R=\overline D-b\ge0$ and
let $\mu_l>0$ be an established lower bound for $\mu$. The endpoint
floors may first be improved by
\begin{align*}
 A_l&\leftarrow\max\left(A_l,
   \sqrt{\frac{(1-\Theta)l}{1-l}}-\frac{R}{6\mu_lS_*}\right),\\
 B_l&\leftarrow\max\left(B_l,
   \sqrt{\frac{(1-\Theta)(1-u)}u}-\frac{R}{6\mu_lS_*}\right),
\end{align*}
using the old $S_*$ in the denominator. Recompute $g,h_{n,\pm}$ from
these proved floors. Combining \eqref{loc:affine-resource} with
$a_0+a_1\ge s_0/O_0+s_1/O_1$ gives
\[
 \lambda_n\{(1-\alpha)C_0(p)+\alpha C_1(p)\}\le R,
 \quad C_0(p)=h_{n,-}+\frac{g}{n(1-p)O_0},\quad
 C_1(p)=h_{n,+}+\frac{g}{npO_1}.
\]
For fixed positive caps, $C_0$ increases and $C_1$ decreases in $p$.
Thus the uniform allocation plane is
\begin{equation}\label{loc:uniform-allocation}
 \lambda_n\{C_0(l)+[C_1(u)-C_0(l)]\alpha\}\le R.
\end{equation}
When the separately checked slope $C_1(u)-C_0(l)$ is nonnegative,
$\lambda_n\le R/C_0(l)$. More generally, positivity of both
endpoint coefficients always gives
\begin{equation}\label{loc:two-endpoint-noise}
 \lambda_n\le\frac{R}{\min\{C_0(l),C_1(u)\}}.
\end{equation}
This consequence permits either ordering of the signed noise prices.
It also contains the scalar harmonic bound. Indeed, for
$B_* =\max(O_0,O_1)\max(u,1-l)$ and
$h_* =\min(h_{n,-},h_{n,+})$,
\begin{equation}\label{loc:harmonic-consequence}
 \min\{C_0(l),C_1(u)\}\ge h_*+\frac{g}{nB_*},\qquad
 \lambda_n\le\frac{R}{h_*+g/(nB_*)}.
\end{equation}
This follows from $(1-l)O_0\le B_*$ and $uO_1\le B_*$.
It improves $\theta$ through
\eqref{loc:canonical-identities}, and hence improves the next endpoint,
outer-cap and price estimates. A prescribed finite number of these
steps proves the strengthened rows by induction; each step uses the
previous variance bound. It assumes neither a fixed point nor a
special value of the unknown allocation $\alpha$.

The saved floor-$25$ enclosures are interpreted as
\begin{equation}\label{loc:weighted-transfer}
 \mathcal R_n\le r_{25},\qquad nD\le25D_{25}.
\end{equation}
This is a fixed sequence of inequalities for the same actual-$n$ law.
Appendix~\ref{ut:section} gives the complete low/high-frequency
candidates, their index gates and the finite stage induction.
The first unseeded raw bound licenses removal of the Gaussian weight
before its defect output is used. Every subsequent step first
transfers the incoming physical bounds, acquires the needed support
at the actual integer, and then restores the weighted defect bound.
Neither extremality at $25$ nor monotonicity of an adaptively rounded
program is assumed. Physical deductions may use the weaker $D\le D_{25}$.

For the exported plane set $r=25/n$, $K=\min(K_-,K_+)$,
$h_{25}=K(1-\Theta)/(25v_*)$ and $R=D_{25}-b$.
If $R>0$, write $d=R/g$ and $L=R/h_{25}$. Then
the resulting scaled plane is
\begin{equation}\label{loc:scaled-plane}
 \frac{a_0+a_1}{rd}+\frac{\lambda_n}{L}
 \le1-(n/25-1)\frac bR\le1.
\end{equation}
The joint noise intercept $L$ stays fixed. Appendix~\ref{ut:section}
derives this plane, its zero-residual boundary and the weighted
specialization of \eqref{loc:harmonic-consequence}.

Finally, new physical support can improve these same parameters.
From an established $\theta\le t<1$, choose
$0<\ell_l\le\sqrt{(1-t)/v_*}$. Then
\[
 O_0\le R_-/\ell_l-l,\qquad
 O_1\le R_+/\ell_l-(1-u).
\]
Intersect with previous caps and apply the same endpoint elimination;
\eqref{loc:canonical-identities} then updates the variance and mean
bounds. Each step consumes the preceding bound. A scalar consequence
does not replace the joint intercept $L$ in \eqref{loc:scaled-plane}.

Consequently a finite acquisition must cover the entire initial moment
and probability domain by closed cells; prove the support, Fourier and
coefficient inequalities on each cell; and either exclude it by a strict
contradiction or export the simultaneous quantities above. Its surviving
rows are necessary conditions on the original global extremizer.
For the surviving middle-band profiles, the acquired contact radii
satisfy $H_-,H_+<2$, the original $\Lambda<1/2$, and
$p\in[.32,.43]$. Since $\lambda_n\le\Lambda$ and $n\ge26$,
$\sqrt{pq+\lambda_n/n}<3/5$. Hence the raw threshold belongs to
\[
 np+H/\ell\in(np-6/5,np+6/5)\subset(0,n).
\]
At this raw threshold the comparison uses
$np+H/\ell=k+t$ with $|t|\le1/2$, including both choices at a tie.
Every nearest integer lies in $0\le k\le n$; all such integers are
covered or excluded by the acquired contact bounds.
This supplies the continuous meaning of the finite
comparison inputs, independently of their file labels.

% Substantive retained derivation: work/berry-esseen-low26-30-20261026/source/
% cutoff25_20261011/new_root/initial25_b1_exact_v1/RETAINED_ANALYTIC_FRAMEWORK.tex,
% contact/support, identity_fourier, resource_enclosure, and joint-gap sections.
% The fixed-floor interpretation and exact beta0 additions are in source/analytic/
% UNIFORM25_CANTELLI_PRODUCTGAP_REVIEW.txt and
% UNIFORM25_BETA0_SOURCE_AND_COMPOSITION_REVIEW.txt.
% Scaled inward plane: source/original/new_audit/
% SCALED_INWARD_PHYSICAL_DERIVATION_V1.txt.

\section{The five finite-index steps}\label{fc:section}

We now compose the comparisons at the actual integers
$n=26,\ldots,30$. The predecessor in each step is the unrestricted
bound $C_{n-1}\le\CE$. All parameters in a single step describe
one extremizer supplied by Lemma~\ref{gi:extremizer}.
In particular, neither its third moment nor its noise is chosen
independently for different count layers.
Throughout this section $R(t)$ denotes \eqref{pp:actual-residual}
with the fixed target $c_n=\CE/\sqrt n$. Computations can use
a smaller directed rational coefficient: an upper bound for that
larger residual is also an upper bound for this $R(t)$.

\subsection{Localization and the domain of the calculation}
The supporting localization proof gives three closed moment ranges.
For $\rho\ge6/5$, a support-free Prawitz comparison at each actual
integer gives $R_n(P)<.4097<\CE$; the large-moment ray is included.
For $1\le\rho\le11/10$, the contact and resource reductions,
with the full predecessor bound, give complete conditional count
covers. The original covers across the five integers consist of $79$ probability
intervals, $2289$ count roots and $74170$ terminal inequalities.
Every other probability interval is excluded already during
localization. These counts describe the coverage, not a replacement
for the inequalities proving it. The present composition uses those
ordinary count proofs on $72$ intervals.
Proposition~\ref{fc:lowprofiles} covers the remaining seven.

It remains to handle $11/10\le\rho\le6/5$.
At each actual integer, localization covers the $100$ original
closed intervals $[i/100,(i+1)/100]$, $0\le i<100$.
It excludes $89$ intervals and leaves exactly
\[
 I_i=[i/100,(i+1)/100],\qquad 32\le i\le42.
\]
The four outer intervals $I_{32},I_{33},I_{41},I_{42}$ have
complete ordinary count covers. At $n=26$, the same is true of
$I_{40}$. The new comparisons therefore concern $34$ parameter
profiles: $I_{34},\ldots,I_{39}$ at $26$, and
$I_{34},\ldots,I_{40}$ at each of $27,28,29,30$.
Their scalar caps are listed in Appendix~\ref{app:parameters}.

A profile supplies inward and outward support caps
$\Delta,O_0,O_1$, a noise cap $\Lambda$, and the necessary joint
inequalities
\begin{equation}\label{fc:resources}
 a_0+a_1\le\Delta(1-\lambda/\Lambda),\qquad
g(a_0+a_1)+\lambda
       \{(1-\alpha)\mathfrak h_-+\alpha\mathfrak h_+\}
       \le\mathfrak R.
\end{equation}
The second inequality is the consequence of
\eqref{loc:affine-resource} with $D\le\overline D$ and
$\mathfrak R=\overline D-b$. It can strengthen the first cap and supplies
the signed plane \eqref{mc:signed-plane}. Its coefficients come
from the actual-index resource proof. They are held fixed on
their stated probability interval, not fitted to the count being
estimated. The contact bounds also give
\begin{equation}\label{fc:contact}
 -H_-\sqrt{pq+\lambda/n}
 \le k+t-np\le H_+\sqrt{pq+\lambda/n},\qquad |t|\le1/2.
\end{equation}
This leaves only three neighboring counts in each profile.
The support caps, both resource bounds and the actual-moment floor
are used simultaneously, with no finite-support assumption on the
original conditional laws.

\subsection{Bounding a whole parameter region}
For a current-count price, let $L_*$ be the acquired actual-noise cap.
The seven original closed noise intervals have endpoints
\[
 0,\quad\Lambda/8,\quad\Lambda/4,\quad5\Lambda/16,\quad
 3\Lambda/8,\quad7\Lambda/16,\quad\Lambda/2,\quad L_*;
\]
the four allocation intervals have endpoints $0,1/4,1/2,3/4,1$.
These endpoints are increasing in every profile.
Their Cartesian product gives $28$ fixed reference laws per
current-count family. Each is valid for its entire closed region
by Lemma~\ref{mc:two-moment}; it is not a sample at its midpoint.
The signed plane and the inward cap restrict the admissible
conditional variances before the reference law is formed.
The noise correction in Proposition~\ref{ce:fixed-box} preserves
the difference between that reference law and the actual noise.
Probability faces and phase faces extend to the whole region
by the corresponding lower-curvature bounds. A tilted phase
face at the Gaussian inflection is used when needed.

For a reflected-square price, Proposition~\ref{ce:coherent}
keeps the actual noise and allocation until the two derivative
gates have been proved. The resulting endpoint envelope uses
the same moment floor as the original law. A positive-atom
activation is covered by the weak derivatives
\eqref{pp:activation-curvature}. Complementary prices may then be
combined through the finite lower envelope of Lemma~\ref{pp:mixture},
using their directed lower target coefficients. Only explicit reference prices are
differentiated in these steps.

All parameter endpoints, chord coefficients and interval enclosures
are rational. The square roots in the finite prices, and the
normal CDF and density, are enclosed outwards. A stationary
width is accepted only when it belongs to its asserted active
set. Curvature bounds include activation boundaries. The finite
maximum of every interpolating quadratic includes its interior
vertex when present. Thus a negative upper bound applies to
every point of the closed domain, including all shared faces.
The full count roots and the unions of their closed phase
intervals are retained when these regional bounds are composed.

\subsection{Examples of the finite comparisons}
An illustrative replacement occurs at $n=26$, $p\in I_{35}$,
count $9$. The current-count price with the least legal slope
is strict on $[-1/8,9/32]$ except for the last phase strip of
the first noise interval. On that first noise interval,
$[0,\Lambda/8]$, retain the current price through $t=1/4$ and
use a coherent count-$10$ reflected-square price of width $9/8$
on $[1/4,9/32]$. Its allocation and noise derivative margins
are respectively at least $.120534962164$ and $.072056354791$,
and its variance derivative is at least $.759364977755$.
Its residual upper bound is $-.018464965470$ for its directed
lower target coefficient. The other original noise intervals
retain the current price on the whole strip.
The resulting uniform upper bound on $[-1/8,9/32]$ is
\begin{equation}\label{fc:n26-simple}
 R(t)\le-.000067780727.
\end{equation}
The adjacent negative phases have a full-square or current-count
bound, and $[9/32,1/2]$ has the full-square count-$10$ bound.
Consequently this entire count layer uses only second-moment
comparisons. The remaining count-$10$ regions in $I_{36},I_{37}$
likewise admit current-count comparisons with the rational slope
multipliers $1,5/4,3/2$. No fourth-moment price or optimized
deterministic-group slope is needed by these analytic profile comparisons.

A second completion illustrates why it is useful to distinguish
the actual CDF from its price. At $n=30$, $p\in I_{37}$,
the count-$10$ contact interval begins at the exact outward
bound
\[
 a_{37}=249467754491/10^{12}<1/4.
\]
An endpoint estimate at $1/4$ extends over $[a_{37},1/4]$
by Lemma~\ref{pp:cdf-bridge}. The comparison used to prove that
endpoint estimate need not remain valid on the short interval.
Monotonicity of the original CDF is sufficient.

\begin{proposition}[Seven lower-moment profiles]\label{fc:lowprofiles}
For each row below, retain the original actual-$n$ support and resource
acquisition and the unrestricted predecessor $C_{n-1}\le\CE$.
The whole profile $p\in I_i$, $1\le\rho\le11/10$ has conditional
count residual bounded by the last column. The joint intercept
$\Lambda$ is retained; $L_*$ is a necessary ceiling on the actual noise.
\begin{center}\small\setlength{\tabcolsep}{3pt}
\begin{tabular}{ccrrrr}\toprule
$n$ & $I_i$ & $\Delta$ & $\Lambda$ & $L_*$ & residual upper\\\midrule
26 & $I_{39}$ & .05769 & .30708 & .17375152 & $-.000098710877$\\
26 & $I_{40}$ & .057809 & .301161 & .17805629 & $-.000005610547$\\
26 & $I_{41}$ & .054236 & .276951 & .162619 & $-.000049410586$\\
27 & $I_{41}$ & .052227 & .276951 & .16129388 & $-.000016889$\\
28 & $I_{41}$ & .050362 & .276951 & .16009718 & $-.000027176644$\\
29 & $I_{41}$ & .048625 & .276951 & .15901083 & $-.000059169468$\\
30 & $I_{41}$ & .047004 & .276951 & .15802002 & $-.000017416596$\\
\bottomrule\end{tabular}\end{center}
\end{proposition}
\begin{proof}
Appendix~\ref{app:finite-details} gives all seven finite plans and their
closed count/phase unions. Each uses the seven original noise bands and
four allocation quarters. The signed resource plane first gives $L_*$;
current-count and reflected-square prices then cover every surviving
contact interval, with the bounds in the table.

Three features require the full comparison rules: reversed signed
prices use \eqref{loc:two-endpoint-noise}; the count-$12$ reference at
$n=29$ crosses $np$ and uses both branches of
\eqref{ce:variance-endpoint-max}; and the early $n=30$ square has
$y_{\max}>0$, covered by the $y_{\max}<b$ bound in Section~\ref{pp:section}.
The remaining steps use the same endpoint, curvature and threshold-shift
identities, with each row's actual integer and acquired parameters.
\end{proof}

\subsection{The finite implications}


The complete middle-moment profiles at $27$--$30$ give the following outward
upper bounds on $R(t)$, uniformly over their probability interval,
actual noise, allocation, all surviving counts and contact phases.
For readability, the table reports $10^6$ times these bounds,
rounded upwards to three decimal places.
\begin{center}
\begin{tabular}{crrrr}
\toprule
interval & $n=27$ & $n=28$ & $n=29$ & $n=30$\\
\midrule
$I_{34}$ & $-320.741$ & $-263.994$ & $-321.528$ & $-241.967$\\
$I_{35}$ & $-297.200$ & $-134.578$ & $-287.407$ & $-436.298$\\
$I_{36}$ & $-294.861$ & $-12.020$  & $-96.700$  & $-156.955$\\
$I_{37}$ & $-145.182$ & $-12.884$  & $-1.810$   & $-63.857$\\
$I_{38}$ & $-70.159$  & $-56.159$  & $-79.848$  & $-22.937$\\
$I_{39}$ & $-35.366$  & $-4.435$   & $-585.317$ & $-75.052$\\
$I_{40}$ & $-1160.650$& $-723.488$ & $-713.730$ & $-987.882$\\
\bottomrule
\end{tabular}
\end{center}
The six analytic profiles at $26$ are also strictly negative;
\eqref{fc:n26-simple} describes its only former fourth-moment
region. The full results include the exact resource caps, count
domains, price choices, derivative gates and outward rational
bounds, rather than only these decimal margins.

\begin{proposition}[The five actual-index implications]\label{fc:five}
For each $n=26,27,28,29,30$,
\[
 C_{n-1}\le\CE\quad\Longrightarrow\quad C_n\le\CE.
\]
\end{proposition}
\begin{proof}
If $C_n>\CE$, Lemma~\ref{gi:extremizer} supplies a positive
closed extremizer. The full predecessor bound makes its
allowance \eqref{gi:predecessor} zero, so the actual-index
localization and resource inequalities apply to this same law.
The support-free Fourier estimate excludes $\rho\ge6/5$.
The complete lower-moment count covers exclude
$1\le\rho\le11/10$.
In the remaining closed band, localization either gives an
immediate contradiction or places the law in one of the eleven
intervals above. The ordinary count covers and the $34$ analytic
profiles exclude them all. Every surviving count and its entire
contact interval has a strict negative bound, contradicting the
positive extremal score. The three closed moment ranges exhaust
$\rho\ge1$.
\end{proof}

\begin{theorem}[The iid Esseen-constant inequality]\label{fc:main}
For every standardized real distribution $P$ with
$\rho(P)=\E|X|^3<\infty$ and every positive integer $n$,
\[
 \sqrt n\,\Delta_n(P)\le
 \frac{3+\sqrt{10}}{6\sqrt{2\pi}}\,\rho(P).
\]
The constant is optimal over this class and all $n$.
\end{theorem}
\begin{proof}
Input~(I) of Proposition~\ref{gi:inputs} supplies the unrestricted
bound at each index through $25$. Apply Proposition~\ref{fc:five}
in order, discharging each global predecessor, to obtain the
bounds through $30$. In particular $C_{30}<\CE+1/100$, so
input~(II) gives $31$--$46$. Proposition~\ref{gi:tail} then gives
every $n\ge47$ from its two sharp predecessors and the uniform
local exclusion (III).
These four ranges exhaust all positive integers.

For the standardized Bernoulli law with
$p_*=(4-\sqrt{10})/2$, the classical binomial asymptotic gives
\[
 \lim_{n\to\infty}\frac{\sqrt n\,\Delta_n(P_*)}{\rho(P_*)}=\CE;
\]
see Schulz \cite[Remarks~4.1 and 4.2]{Schulz2016}.
Thus a smaller uniform constant is impossible.
\end{proof}

\section{Scope of the supporting material}\label{fc:supporting}
The finite comparisons in Sections~\ref{mc:section}--\ref{ce:section}
are analytic inequalities for arbitrary real conditional laws.
The finite calculations certify their use on the closed parameter
domains of Section~\ref{fc:section}. The supporting material also
contains the continuous localization and the proofs of inputs
(I)--(III). These are essential parts of the all-index argument,
not hypotheses inferred from the numerical table.

The moment-defect and localization framework, the original
subdivisions, and the imported count and Fourier arguments are
due to Rodziewicz and Ko\l{}odziejek \cite{RK2026}.
The role of the present comparison method is to close the remaining
general-law regions at the five actual integers and compose the
supporting implications in a well-founded order.

This working version includes the companion mathematical sources,
exact acquired parameters, collected frozen arithmetic inputs and
archived calculation outputs. The accompanying programs locate these
materials and run their mathematical readers from the collected
directories. Fresh numerical regeneration uses the retained producer
entrypoints and their stated runtime requirements. The collection is
local at this stage; no permanent public archive identifier is assigned
in this draft.


\section*{AI assistance}
OpenAI Codex was used to develop and simplify mathematical arguments,
implement exact-arithmetic calculations, and draft this manuscript.

\appendix
\section{A fixed-floor Fourier and resource transfer}\label{ut:section}

This appendix proves the interpretation of the fixed floor-$25$
acquisition used in \eqref{loc:weighted-transfer}. The floor is an
evaluation parameter in a collection of inequalities. It is not another
sample size at which the current law is assumed to be extremal.
The moment-defect identity and the underlying enclosure method retain
their attribution to Rodziewicz--Ko\l{}odziejek~\cite{RK2026}.

Fix $N=25$, an integer $n\in\{26,\ldots,30\}$, and a standardized law
$P$ attaining the unrestricted $C_n=C\ge\CE$ at a positive closed
threshold $H$. Assume $C_{n-1}\le\CE$ and
$\rho\in[1,11/10]$ or $\rho\in[11/10,6/5]$.
Use $D$, $\mathcal L$ and $\mathcal R_n$ from
\eqref{loc:conserved}--\eqref{loc:raw}, all for this actual pair $(P,H)$.
Write $\mathcal R_n=\mathcal R$ for the right side of
the conserved identity \eqref{loc:conserved} at this same law.
We require an actual-index score cap $C\le c_{\rm in}$ with
$c_{\rm in}\le434477/10^6$; the Gaussian replacement bound
$C<.422$ under the sharp predecessor is sufficient.
The support and moments used below must have been acquired for this
same actual-$n$ pair. No assertion about $C_{25}(P)$ is required.

\subsection{The complete low-frequency estimate}

Put $f(t)=\E e^{itX}$, $\gamma=\E X^3$, and
\begin{align*}
 V_n(t)&=f(t)^n-e^{-nt^2/2}(1-in\gamma t^3/6),\\
 I_0(n,H)&=\int_0^{1/2}\frac{|\Im(e^{-itH}V_n(t))|}{t^3}\,dt,\\
 I_\infty(n)&=\int_{1/2}^{\infty}\frac{|V_n(t)|}{t^3}\,dt.
\end{align*}
By \eqref{loc:fourier-remainder},
\begin{equation}\label{ut:split}
 |E_n|\le\frac2\pi\{\sqrt n I_0(n,H)+\sqrt n I_\infty(n)\}.
\end{equation}
Both integrals are retained, including the entire Gaussian comparator
on $[1/2,\infty)$.

\begin{lemma}[Low-frequency candidates]\label{ut:low}
Suppose $n>2$, $M_4\le M\le111/10$, $|\gamma|\le g\le37/20$,
and $M_6\le\overline M_6$. For the actual coefficients
\[
 a=(3-M_4)/24,\qquad z_5=(10\gamma-\E X^5)/120,
 \qquad c_5=a\gamma/6,
\]
let $a_*\ge|a|$ and $Z\ge|z_5|$. Define, for real $j>2$,
\begin{gather*}
 A_j=\frac{j-1}{2j},\quad B_j=\frac{45-M}{96}(1-2/j),
 \quad F_j=(j/(j-1))^{3/2},\quad I_u\ge\sqrt{\pi/2},\\
 E_6=\frac{\overline M_6}{720}+\frac{g^2}{72}
                  +\frac{(1/2+g^2/144)^3}{5},\quad
 T_6=a_*+Z/2+E_6/4,\quad
 L_2(j)=\frac{E_6}{2A_j^2}+\frac{T_6^2}{2B_j^3}.
\end{gather*}
Each of the following complete triples satisfies
\begin{equation}\label{ut:low-bound}
 \sqrt n I_0(n,H)\le L_0(n)+L_H(n)|H|+L_2(n)/\sqrt n.
\end{equation}
The triangle triple is
\begin{align}\label{ut:triangle}
 L_0(j)&=I_uF_j(Z+a_*g/2)+251g^3/17280,\notag\\
 L_H(j)&=I_uF_j a_*+251g^2/4800.
\end{align}
For the cancellation triple, let a closed rectangle contain the actual
$(z_5,c_5)$, and let $J_*\ge0$ bound
\[
 J(z,c)=\int_0^\infty e^{-x^2/2}|zx^2-cx^4|\,dx
\]
at every corner of that rectangle. Then one may take
\begin{align}\label{ut:paired}
 L_0(j)&=F_jJ_*+\frac{F_jI_u}{j-1}
       (35a_*g^3/48+35Zg^2/24)+251g^3/17280,\notag\\
 L_H(j)&=F_jI_u\{a_*+5Zg/[2(j-1)]\}+251g^2/4800.
\end{align}
In either case, all three coefficients and $L_2(j)/\sqrt j$ are
nonincreasing for $j\ge N$ when their physical inputs are fixed.
\end{lemma}
\begin{proof}
Set $h(t)=e^{-t^2/2-i\gamma t^3/6}$. Taylor expansion gives
\[
 f(t)-h(t)=-at^4-iz_5t^5+r_6(t),\qquad
 |r_6(t)|\le E_6t^6\quad(0\le t\le1/2).
\]
Indeed, for $u=-t^2/2-i\gamma t^3/6$,
$|u|\le t^2(1/2+g^2/144)<1/6$;
the omitted exponential terms cost at most $|u|^3/5$, and its
degree-six quadratic term costs $g^2t^6/72$.
Symmetrization and $\cos y\le1-y^2/2+y^4/24$ give
\[
 |f(t)|,|h(t)|\le e^{-(45-M)t^2/96}.
\]
Taylor expansion of the polynomial $w^n$ along the joining segment
therefore yields
\[
 f(t)^n-h(t)^n=n(f(t)-h(t))h(t)^{n-1}+R_2(t),\qquad
 |R_2(t)|\le\tfrac12n^2T_6^2t^8e^{-nB_nt^2}.
\]
Integrating $nr_6h^{n-1}$ and $R_2$, after division by $t^3$,
costs at most $L_2(n)/n$.
For $u=n\gamma t^3/6$, rotation by $Ht$ bounds the remaining
Gaussian comparator's imaginary part by
$e^{-nt^2/2}(|u|^3/6+|H|t u^2/2)$.
The Gaussian even moments give its integral at most
$n^{-1/2}(251g^3/17280+251g^2|H|/4800)$.

For the leading quartic and fifth terms put $b=n-1$ and
$v=b\gamma t^3/6$. The bounds
$|\sin(Ht+v)|\le|H|t+bg t^3/6$ and $|\cos|\le1$
give \eqref{ut:triangle}.
Alternatively, the leading imaginary part after division by $t^3$
is bounded, before multiplication by $ne^{-bt^2/2}$, by
\[
 |z_5t^2-atv|+a_*t|\sin v-v|+Zt^2|\cos v-1|
                 +|H|t\{a_*t+Zt^2|\sin v|\}.
\]
The first integral is $nJ(z_5,c_5)/b^{3/2}$.
The norm $J$ is convex in $(z,c)$, so its rectangle maximum is
bounded by the largest of its corner bounds. Integrating the other
four terms with the global sine and cosine remainders gives
\eqref{ut:paired}. These arguments use the original $H$, without
rescaling it to a floor-$N$ threshold.

Finally,
\[
 \frac{F_j'}{F_j}=-\frac3{2j(j-1)},\quad
 A_j'=\frac1{2j^2}>0,\quad B_j'=\frac{45-M}{48j^2}>0.
\]
Thus $F_j$, $F_j/(j-1)$, $A_j^{-2}$, $B_j^{-3}$ and
$j^{-1/2}$ decrease. Every multiplier of these factors in
\eqref{ut:triangle}--\eqref{ut:paired} is nonnegative, including
$J_*$. This proves each claimed coefficient comparison.
\end{proof}

Here are the coefficient choices actually needed in these candidates.
They are physical moment inequalities, independent of the sample size.
For the triangle choice, assume $|X|\le R\le6$ and $M_5\le\eta$.
One may use $a_* =\max(2,M-3)/24$, $\overline M_6=R\eta$, and
$120Z=\min\{\eta+10g,Z_6,Z_r\}$, where
\[
 Z_6=\begin{cases}
 6(M-1)+(10-2M)g-6g^2+g^3,
       &M\le5,\ 10-2M-12g+3g^2\ge0,\\
 6(M-1)+(10-2M)^2/[4(6-g)],&M\le5\text{ otherwise},\\
 6(M-1)+(2M-10)^2/24,&M>5,
 \end{cases}
\]
and, with $r_0=\max(R,2g)>g$, put
\[
 t_*=\begin{cases}
 \min(g,(10-2M)/(2(r_0-g))),&M\le5,\\
 \min(g,(2M-10)/(2r_0)),&M>5.
 \end{cases}
\]
Then
\[
 Z_r=\begin{cases}
 r_0(M-1)+(10-2M)t_*-(r_0-g)t_*^2,&M\le5,\\
 r_0(M-1)+(2M-10)t_*-r_0t_*^2,&M>5.
 \end{cases}
\]
To verify them, reflect only the auxiliary Fourier integral so that
$t=\gamma=|\gamma|\ge0$, and set $Q=X^2-tX-1$.
Direct multiplication gives
\begin{equation}\label{ut:pearson}
\begin{aligned}
 \E Q^2&=M_4-1-t^2,\\
 \E X^5&=t^3+2t+\E[Q^2(X+2t)],\\
 M_6&=1+3t^2+t^4+\E[Q^2\{(X+t)^2+2t^2+2\}].
\end{aligned}
\end{equation}
For an enclosing radius $r\ge2g$, the two signs of
$\E X^5-10t$ are bounded by
$r(M-1)\pm(10-2M)t-rt^2\pm t^3$.
Use $t^3\le gt^2$, or discard the negative cubic, and maximize the
resulting concave quadratic on $[0,g]$. At $r=6$, the increasing
endpoint alternative gives the first line of $Z_6$.

The resource implementation also uses the following general signed
rectangle, which does not require $M<3$ or $R>2g$.
Suppose $0\le g_l\le|\gamma|\le g$, $M\ge1+g_l^2$, and $R\ge1$.
From \eqref{ut:pearson}, for every feasible $t\in[g_l,g]$,
\begin{gather*}
 L(t)\le120z_5\le U(t),\\
 L(t)=-R(M-1)+(10-2M)t+Rt^2+t^3,\\
 U(t)=\begin{cases}
 R(M-1)+(10-2M)t-Rt^2+t^3,&t\le R/2,\\
 8t-t^3,&t\ge R/2.
 \end{cases}
\end{gather*}
The second line for $U$ uses $\E Q^2\ge0$ when $R-2t<0$;
using $M-1-t^2$ there would have the wrong direction.
Choose outward bounds $L_*\le\min L$ and $U_*\ge\max U$ on the
indicated closed intervals, and set
\begin{gather*}
 z_l=\max(-\eta-10g,L_*)/120,\quad
 z_u=\min(\eta+10g,U_*)/120,\quad Z=\max(|z_l|,|z_u|),\\
 a_* =\max(|3-M|,|2-g_l^2|)/24,\\
 c_l=\frac{3-M}{144}
       \begin{cases}g_l,&M\le3,\\g,&M>3,\end{cases}
 \qquad c_u=\max_{g_l\le t\le g}(2t-t^3)/144,\\
 \overline M_6=\min\{R\eta,
  1+3g^2+g^4+(R^2+2Rg+3g^2+2)(M-1-g_l^2)\}.
\end{gather*}
Check $z_l\le z_u$, $c_l\le c_u$ before using the rectangle.
The cubic extrema can be enclosed by the Bernstein identity on a
finite closed subdivision; this is an enclosure of the entire
interval, rather than evaluation at finitely many sample points.
The displayed $a_*$ uses $1+g_l^2\le M_4\le M$.
The last line follows from the final identity in \eqref{ut:pearson},
with all substituted factors nonnegative. These prove the hypotheses
of Lemma~\ref{ut:low} for the general signed candidate used by both
resource branches. The restricted $M<3$, $R>2g$ formula is not needed
to license a call outside its own domain.

The resource branch enables this signed candidate only on an already
acquired interval $0<l\le p\le u<1/2$, with $\theta\le\Theta<1$
and $A_-\ge A_l>0$, $A_+\ge B_l>0$.
Write $\rho_B(p)=\{1-2p(1-p)\}/\sqrt{p(1-p)}$ and set
\begin{gather*}
 \rho_* =\max\left\{1,\rho_l,\rho_B(u)
       -\max(0,3\rho_B(l)/2-3\min(A_l,B_l))\Theta\right\},\\
 g_l=\max\left\{0,\sigma_l\rho_*-3-d_D,
 \frac{(1-\Theta)^{3/2}(1-2u)}{\sqrt{u(1-u)}}
                -(R+2\sqrt{u/(1-u)})\Theta\right\}.
\end{gather*}
Here $D\le d_D$ is an incoming bound. These are genuine physical
lower bounds for $\rho$ and $|\gamma|$, respectively.
For the first, conditional expansion gives
$\rho\ge(1-\theta)^{3/2}\rho_B(p)+3\min(A_l,B_l)\theta$;
use $(1-\theta)^{3/2}\ge1-3\theta/2$ and the monotonicity of
$\rho_B$ on $(0,1/2)$.
For the second, $\gamma\ge\sqrt{10}\rho-3-D$ follows from
$J\le\rho+1$, whereas conditional expansion about the two means gives
\[
 \gamma\ge\frac{(1-\Theta)^{3/2}(1-2u)}{\sqrt{u(1-u)}}
                     -(R+2\sqrt{u/(1-u)})\Theta.
\]
In the latter calculation the positive-class residual is nonnegative
and the negative-class residual has magnitude at most the displayed
coefficient times $\theta$.
The gates are $\rho_*\le\rho_u$, $0\le g_l\le g\le37/20$,
$1\le M\le111/10$, $R\ge1$ and $M-1-g_l^2\ge0$.
The same $\rho_*$ is kept in the paired rectangle, score corners,
phase callback and exports. On another probability half, or if this
candidate is ineligible, the unpaired candidate remains available;
ineligibility alone does not exclude a law.

For completeness, the smaller Gaussian corner prices used in this
rectangle also have an index-free proof. If $zc\le0$, then
$J(z,c)=\sqrt{\pi/2}(|z|+3|c|)$.
If $zc>0$, replace $(z,c)$ simultaneously by $(|z|,|c|)$ and put
$x=z/c>0$. Splitting the integral at $\sqrt x$ gives
\begin{equation}\label{ut:mills-norm}
\begin{aligned}
 J(z,c)&=\sqrt{\pi/2}(z-3c)+\sqrt{zc}e^{-x/2}
       \left(6+\frac{2(3-x)}x\sqrt x\,\mathcal M(\sqrt x)\right),\\
 \mathcal M(r)&=\int_0^\infty e^{-rt-t^2/2}dt.
\end{aligned}
\end{equation}
For $I_j(r)=\int_0^\infty t^je^{-rt-t^2/2}dt$, integration by parts
gives $I_0=1/(r+I_1/I_0)$ and
$I_j/I_{j-1}=j/(r+I_{j+1}/I_j)$.
Truncating this finite recursion with a zero last ratio gives an
alternating upper bound at even depth and lower bound at odd depth.
In rational form, set $d_{k+1}=1$,
$d_j=1+j/(xd_{j+1})$ for $j=k,\ldots,1$.
In \eqref{ut:mills-norm} replace
$\sqrt x\mathcal M(\sqrt x)$ by $d_1^{-1}$ with $k=11$ if
$x\ge3$, and $k=12$ if $x<3$.
The sign of $3-x$ gives the required upper bound in both cases.
The remaining multiplier is positive; for $x\ge3$ it exceeds
$6-2(x-3)/x>4$.
A negative coefficient $z-3c$ uses a lower enclosure of
$\sqrt{\pi/2}$ and a positive coefficient uses an upper enclosure.
The minimum with
\[
 \sqrt{\pi/2}|z-3c|+6\sqrt{zc}e^{-x/2}
\]
is also a valid corner bound. These are the two corner alternatives
in the saved candidates. Auxiliary reflection here leaves the original
signed $\gamma$ in the contact identity unchanged.

\subsection{The complete high-frequency estimate}

Let $X'$ be an independent copy of this same $X$. A candidate consists
of fixed numbers $M,a,m,d,w$ and an actual event $E$ such that
\begin{equation}\label{ut:event-input}
\begin{gathered}
 M_4\le M\le111/10,\qquad\PP(E)\ge a>0,\\
 |X-X'|\ge d>0\text{ on }E,\qquad
 \E[(X-X')^2\mid E]\le m,\\
 0<d\le12,\qquad0<w\le1/4.
\end{gathered}
\end{equation}
Here $a\le1$ and $m>0$. Put $b=9/10$,
\begin{gather*}
 h=333/106-(22/7)^3w^2/6>0,\quad
 \kappa_E=a(23/48)^2d^2,\quad
 \Lambda(x)=\frac{x+2}{x^2+4x+2},\\
 S_3=\sum_{j=1}^{20}j^{-3}+1/800,\quad
 S_5=\sum_{j=1}^{20}(j-w)^{-5}+[4(20-w)^4]^{-1},\\
 c=\tfrac12\left(1-\max_{u\in\{1/4,81/100\}}
                    [1-u+(M+3)u^2/12]\right),\\
 c_0=1/8-(M+3)/384,\qquad s=\{1-(M+3)53/600\}/2.
\end{gather*}
Define
\begin{align}\label{ut:high-six}
 \mathfrak H_j={}&A_j^{\rm high}
 +\frac{m}{4\pi^{3/2}h\sqrt a}
                \left(S_3+\frac{3S_5}{jah^2}\right)
 +\frac{m\sqrt j}{2\pi^2}e^{-jah^2w^2}\notag\\
 &+Z_j+\frac{\sqrt j}{2b^2}e^{-ja(23/48)^2}
 +\sqrt j e^{-j/8}\{2\Lambda(j/8)+\sqrt{M-1}/3\}.
\end{align}
One always may choose
\[
 A_j^{\rm high}=2\sqrt j e^{-jc},\qquad
 Z_j=\frac{e^{-j\kappa_Eb^2}}{2b^4\kappa_E\sqrt j}.
\]
If $s>0$, the first may be replaced by its minimum with
$2\sqrt j e^{-jc_0}\Lambda(js/4)$.
If $j\kappa_Eb^2\ge1/4$, the second may be replaced by its minimum
with $\sqrt j e^{-j\kappa_Eb^2}\Lambda(j\kappa_Eb^2)/(2b^2)$.

\begin{lemma}[High-frequency candidates and their gates]\label{ut:high}
Under \eqref{ut:event-input}, $\sqrt n I_\infty(n)\le\mathfrak H_n$.
A fixed choice of alternatives satisfies
$\mathfrak H_n\le\mathfrak H_N$ for every $n\ge N$ provided
\begin{equation}\label{ut:high-gates}
 2Nc'>1\quad
 (c'\in\{c,ah^2w^2,a(23/48)^2,1/8\}).
\end{equation}
For the improved first term also require $s>0$ and $2Nc_0>1$;
for the improved zero-cell term require $N\kappa_Eb^2\ge1/4$.
An alternative failing its floor gate is omitted, even if it would
become eligible at a larger index.
\end{lemma}
\begin{proof}
On $[1/2,b]$, symmetrization gives
$|f(t)|^2\le1-t^2+(M+3)t^4/12$, proving the first term with rate
$c$. The sharper chord is
\[
 t^2/2-(M+3)t^4/24\ge c_0+s(t^2-1/4).
\]
It is valid since the difference is
$(M+3)(t^2-1/4)(81/100-t^2)/24\ge0$.

For $t\ge b$, let $p_E=\PP(E)$. Symmetrization, $1-u\le e^{-u}$,
and conditional Jensen imply
\[
 |f(t)|^n\le e^{-n\E\sin^2(t(X-X')/2)}
 \le\E[e^{-na\sin^2(t|X-X'|/2)}\mid E].
\]
For each conditional distance $D_*=|X-X'|$, partition the positive
axis into the zero cell and the cells centred at
$2\pi j/D_*$, $j\ge1$. In a nonzero cell write
$t=(2\pi/D_*)(j+v)$. On $|v|\le w$,
$|\sin(\pi v)|\ge h|v|$; on the wings it is at least $hw$.
The paired Taylor bound
\[
 \{(j-v)^{-3}+(j+v)^{-3}\}/2
             \le j^{-3}+6(j-w)^{-5}v^2
\]
and the Gaussian zeroth and second moments prove the second term
of \eqref{ut:high-six}. Its Jacobian is
$D_*^2/(4\pi^2)$, and conditional averaging uses
$\E[D_*^2\mid E]\le m$. Summing the two positive series and
bounding their tails after $20$ by integrals gives $S_3,S_5$.
The wing weight is at most $D_*^2/(2\pi^2)$ and gives the third term.
All extensions of these integrals below $b$ add nonnegative quantities.

In the zero cell, $u=tD_*/2\le\pi/2$.
For $u\le1/2$ use $\sin u\ge(23/24)u$ and $D_*\ge d$.
This gives the zero-cell Gaussian tail $Z_j$.
For $u\ge1/2$ use $\sin u\ge23/48$ and total weight
$1/(2b^2)$, giving the fifth term.
The comparator on the whole interval $[1/2,\infty)$ costs
\[
 \sqrt j\int_{1/2}^\infty e^{-jt^2/2}t^{-3}dt
 +\frac{j^{3/2}|\gamma|}{6}\int_{1/2}^\infty e^{-jt^2/2}dt.
\]
The second integral is at most $2e^{-j/8}/j$, and Pearson gives
$|\gamma|\le\sqrt{M-1}$. This proves the sixth term.

The rational Gaussian-tail estimate used above is, for $x=jk b^2>0$,
\begin{equation}\label{ut:rational-tail}
 \sqrt j\int_b^\infty e^{-jkt^2}t^{-3}dt
 \le\frac{\sqrt j e^{-x}}{2b^2}\Lambda(x).
\end{equation}
After substitution its remaining factor is
$\int_0^\infty u(x+u)^{-1}e^{-u}du$.
For $P(u)=u^2-4u+2$ and $D_x=x^2+4x+2=P(-x)$,
division of $D_x^2-P(u)^2$ by $x+u$, followed by
$\int u^ke^{-u}du=k!$, gives
$\int_0^\infty (x+u)^{-1}e^{-u}du\ge(x+3)/D_x$.
The discarded remainder is $P(u)^2/[(x+u)D_x^2]\ge0$.
Consequently the preceding factor is at most
$1-x(x+3)/D_x=\Lambda(x)$, proving
\eqref{ut:rational-tail}. The simpler zero-cell alternative follows
by replacing $t^{-4}$ by $b^{-4}$ after integrating once.

For the floor comparison,
$(\sqrt j e^{-cj})'=(\sqrt j e^{-cj})(1/(2j)-c)$,
and $j^{-1/2}e^{-cj}$ decreases for every $c>0$.
Also $\Lambda'(x)=-(x^2+4x+6)/(x^2+4x+2)^2<0$.
These prove the assertions for the second, third, fifth and sixth
terms and for either first-term choice under the stated gates.
For the improved zero-cell choice the derivative of
$\sqrt x e^{-x}\Lambda(x)$ has, after multiplication by positive
factors, numerator
\[
 4-10x-22x^2-13x^3-2x^4.
\]
It equals $-11/128$ at $x=1/4$ and decreases for $x>0$.
This proves precisely its additional gate.
\end{proof}

The event inputs transfer without an index assumption. For $0<d^2<2$
and $E=\{(X-X')^2\ge d^2\}$, Cantelli's inequality supplies
\[
 V=2M+2,\quad h_0=2-d^2>0,\quad
 a=h_0^2/(V+h_0^2),\quad m=2+V/h_0.
\]
Indeed $Y=(X-X')^2$ has $\E Y=2$ and
$\operatorname{Var}Y\le V$. Applying the elementary one-sided
second-moment bound to $Y-2$ gives $\PP(Y\ge d^2)\ge a$.
For $p_E=\PP(E)$, Cauchy--Schwarz applied to the covariance gives
$\E[Y\mid E]\le2+\sqrt{V(1-p_E)/p_E}\le2+V/h_0$.
This also licenses the older weaker parameters
$a_U=h_0^2/(V+4)$, $m_U=2(V+4)/h_0^2$:
$a\ge a_U$ and
\[
 m_U-m=\frac{(2-h_0)(V+2h_0+4)}{h_0^2}\ge0.
\]
Every event-dependent term of \eqref{ut:high-six} decreases in $a$
and increases in $m$; its gates are preserved when $a$ increases.
Thus no separate universal-event proof is needed, and old certificates
using the weaker parameters remain valid. This domination holds for
the same $M,d,w$; it does not identify different distance choices.
The selected Cantelli pool uses
$d\in\{1/2,11/20,3/5,13/20,7/10\}$ and $w=1/4$.
Only members satisfying \eqref{ut:high-gates} at $N$ are retained.
Replacing $M$ by a fixed outward ceiling on the $1/1000$ grid
is legitimate as a new fourth-moment cap; it is then held fixed.

After a central gap $(-A_-,A_+)$ is proved, another actual event is
$E=\{XX'<0\}$. With $p=\PP(X>0)$, $\mu=\E X_+$ and
$\beta=\E[X|X|]$, its exact conditional distance is
\[
 \PP(E)=2p(1-p),\qquad
 \E[(X-X')^2\mid E]
 =\frac{1+4\mu^2+\beta}{2p}
  +\frac{1+4\mu^2-\beta}{2(1-p)}.
\]
The expression increases in $\mu^2$, is affine in $\beta$, and is
convex in $p$ for $|\beta|\le1$. Incoming interval endpoints
therefore give $m$; the independent bound $m\le2/a$ may be intersected.
An acquired $D\le d_D<(\sigma_l-1)\rho_l$, where
$1<\sigma_l\le\sqrt{10}$, licenses the central gap. Moreover
\begin{equation}\label{ut:gap-input}
\begin{gathered}
 (A_-+A_+)^2\ge\frac{\rho_l-d_D/(\sigma_l-1)}{\mu},\qquad
 (A_-+A_+)^2\ge\frac{2P_*}{3\mu},\\
 P_* =\max\{(\sigma_l-1)\rho_l,\sigma_l\rho_l-\gamma_u\}-d_D.
\end{gathered}
\end{equation}
The first follows from the gap decomposition's
$T_++T_-\le D/(\sqrt{10}-1)$.
The second follows from
$J\le\rho+2\mu A_-A_+$ and hence $6\mu A_-A_+\ge P_*$.
Use an incoming upper bound for $\mu$ and positive lower roots.
The weaker first bound with $\mu\le1/2$, and independently acquired
$A_-+A_+\ge A_l+B_l$, give the other saved distance choices.
The gates are a positive gap, positive mean denominator and $0<d\le12$;
no assumption $d_D<2$ is needed for \eqref{ut:gap-input}.

A finite minimum of complete high-frequency candidates is valid:
each bounds the same $\sqrt n I_\infty(n)$.
Likewise a finite minimum of complete low-frequency triples is valid.
One may not combine the constant of one low triple with the slope of
another without a separate proof. Rational outward evaluations at
$N$ bound the exact displayed expressions at $N$ and hence at $n$;
no monotonicity of the rounding or exponential evaluator is assumed.

\subsection{Raw and charged candidates at the fixed floor}

Choose one complete low triple and one complete high bound at $N$.
For a fixed rational $\omega_\pi\ge2/\pi$, put
\begin{equation}\label{ut:fourier-floor}
 K_F=\omega_\pi\{L_0(N)+\mathfrak H_N+L_2(N)/\sqrt N\},
 \qquad s_F=\omega_\pi L_H(N).
\end{equation}
Lemmas~\ref{ut:low} and \ref{ut:high} prove
$|E_n|\le K_F+s_F|H|$ for the actual law.
The values $16/25$ and $212/333$ used in the two enclosure families
both satisfy the normalization condition.

Here is a precise description of the coefficient estimates which can
be used with \eqref{ut:fourier-floor}. On an incoming physical cell
assume that the coefficient atlas has proved
\begin{equation}\label{ut:physical-coefficients}
 3\beta J/2+3\rho J_1-3J_2-K_q\le A(\rho),\qquad
 K_1\in[k_l,k_u],\quad U\in[u_l,u_u].
\end{equation}
These are inequalities in the moments and support of $P$; they contain
no sample index. For the initial cover they are proved on the signed
$(\rho,\beta,\gamma,\mu)$ rectangles. For the resource cover they use
the conditional moment and gap formulas of Section~\ref{loc:section}.
The full physical derivations are the sections \emph{Signed moments
with a common first moment}, \emph{Preserving central and outer
integral costs} and \emph{Moment and gap coefficients for a whole
cell} in the continuous companion's \emph{Coefficient atlas for the
enclosure}, located in Appendix~\ref{app:companions}.
The central polynomial is $\rho P_1(A_-,A_+)-P_2(A_-,A_+)$, where
\begin{align*}
 P_1&=B^3/3+AB^2/2-A^3/6,\\
 P_2&=B^4/4+AB^3/3-A^4/12,
\end{align*}
where $(A,B)=(A_-,A_+)$.
Its full-rectangle Bernstein bound is a maximum of affine functions
of $\rho$. The outer cost is bounded by
$3\{(\rho-B_l)_+^2T_++(R-\rho)^2T_-\}/8$.
These and the remaining affine moment terms give the convex
$A(\rho)$ used in the corner reduction. Every signed product with
a possibly negative factor retains the appropriate two endpoints.
Thus \eqref{ut:physical-coefficients}, rather than a file status, is
the required physical-coefficient premise of a candidate.

Suppose previous information gives
$\mathcal R_n\le r$, $nD\le Nd$ and $0<r<2N/5$.
Once removal of the Gaussian weight has been licensed, it follows that
\begin{equation}\label{ut:incoming-envelope}
 D\le d,\qquad L_N\le\varphi\le2/5,\quad
 L_N=\frac{997}{2500}(1-5r/(2N))>0.
\end{equation}
The density statement uses the nonnegative phase term in
\eqref{loc:conserved}, $\CE\rho>2/5$ and $n\ge N$.
For the unseeded first pass use $L_N=0$, without assuming $r$ or $d$.

For a sign $\varepsilon\in\{-1,1\}$ and a prospective score $c$, set
\[
 F_\varepsilon(c)=s_F+
 \max_{k\in\{k_l,k_u\},\,u\in\{u_l,u_u\},\,phi\in\{L_N,2/5\}}
                  \varepsilon(ck+\phi u).
\]
The maximum of the two signs may be used for an unsigned bound.
If $\gamma\ge g_0(\rho)$ is an incoming signed lower bound, with
$g_0$ affine in $\rho$ (including a fixed signed interval endpoint),
the correct phase curvature is
\begin{equation}\label{ut:signed-curvature}
 q(c,\rho)=c\rho+\tfrac12
           \min\{L_N g_0(\rho),(2/5)g_0(\rho)\}.
\end{equation}
In particular $g_0(\rho)=3.16\rho-3-d$ is valid since
$J\le\rho+1$. The alternative general-envelope curvature is
\[
 q_{\rm gen}(c,\rho)=c\rho+L_Ng_0(\rho)/2
                    -(2/5-L_N)(d-.16)_+/2.
\]
It is no larger than \eqref{ut:signed-curvature}, because
$(-g_0(\rho))_+\le(d-.16)_+$ for $\rho\ge1$.
Either curvature must be proved strictly positive before division.
Using $L_Ng_0/2$ alone when $g_0<0$ would not be licensed.

For the actual $H$ with sign $\varepsilon$, \eqref{loc:raw} now gives
\begin{equation}\label{ut:raw-candidate}
 \mathcal R_n\le K_F+C A(\rho)+F_\varepsilon(C)|H|
                                      -q(C,\rho)H^2.
\end{equation}
Thus a floor candidate is an envelope for the actual raw threshold.
It is not an estimate of a fictitious $\mathcal R_{25}(P)$.

The same incoming bounds give a conserved-phase price
\begin{equation}\label{ut:rebate}
 e_N(\rho)=\tfrac12(1-5r/(2N))(.4095\rho-d/15)>0.
\end{equation}
Indeed the difference between the weighted and unweighted terms in
\eqref{loc:conserved} is
$n(1-e^{-w})(\CE\rho-\phi(0)D/6)$, $w=H^2/(2n)$.
Use $1-e^{-w}\ge we^{-w}$, the density floor before multiplying by
$\phi(0)$, $D\le d$ and $\phi(0)<2/5$.
Consequently, with $\kappa=997/15000$,
\begin{equation}\label{ut:charged-candidate}
 \kappa nD+n(C-\CE)\rho
 \le K_F+C A(\rho)+F_\varepsilon(C)|H|
                       -\{q(C,\rho)+e_N(\rho)\}H^2.
\end{equation}
The first unseeded pass uses $e_N=0$.

The own-score interval is also conservative at $N$. If a physical
cell forces $D\ge d_l\ge0$ and $\rho\ge\rho_l>0$, then
\begin{equation}\label{ut:own-score}
 C\le .40974+
 \frac{r/N-L_Nd_l/6}{\rho_l}
\end{equation}
whenever the numerator is nonnegative. A negative numerator excludes
the cell. This follows from \eqref{loc:conserved}, since
$r/n-\varphi d_l/6\le r/N-L_Nd_l/6$.
The independent moment envelope and $c_{\rm in}$ may be intersected
with \eqref{ut:own-score}.

For each candidate freeze all physical coefficients before taking the
score and moment corners. For fixed $H$, the right side of
\eqref{ut:raw-candidate} is convex in $c$ and in $\rho$ separately:
$A$ is convex, $F_\varepsilon$ is a maximum of affine functions,
and $q$ is affine in $c$ and concave in $\rho$.
The same is true after adding the affine curvature $e_N$.
Hence the whole physical rectangle is bounded by its corners, first
pointwise in $H$, and then by completing each corner's square.
For $Q>0$, $\sup_{h\ge0}(A+Fh-Qh^2)=A+(F_+)^2/(4Q)$;
using $A+F^2/(4Q)$ is a weaker valid bound.

For a charged candidate split the score interval at $.40974$ and use
\begin{equation}\label{ut:excess-charge}
 n(C-\CE)\rho\ge N(C-.40974)_+\rho.
\end{equation}
On either side of this breakpoint the subtracted right side is
affine in each of $c,\rho$ separately, so the same corner argument
applies. Omitting this breakpoint would not follow from convexity.
If $R_N$ is the maximum raw corner price and $B_N$ the maximum
charged corner price, the resulting conclusions are
\begin{equation}\label{ut:stage-output}
 \mathcal R_n\le R_N,\qquad nD\le B_N/\kappa.
\end{equation}
They remain true after minima with other complete candidate bounds.
The phase roots of these same corner quadratics bound the original
$H$, and the pointwise charged quadratics give
$D\le\max_j(A_j+F_{j,\varepsilon}|H|-Q_jH^2)/(\kappa n)$.
Their normalization is therefore $15000/(997n)$.
For an unpaired export the implementation may instead use an already
proved $|H|\le H_*$ and the scalar weighted defect bound to write
\[
 \kappa nD\le\kappa Nd
 \le\kappa Nd+\epsilon(H_*^2-H^2),\qquad\epsilon>0.
\]
This is a valid pointwise envelope on the acquired phase interval.
Its artificial quadratic term asserts no additional physical curvature.

\subsection{The unseeded step and finite induction}

\begin{proposition}[Fixed-floor interpretation]\label{ut:transfer}
Consider the fixed finite B0 or B1 acquisition schedule at $N=25$.
At each Fourier step retain a complete candidate satisfying
Lemmas~\ref{ut:low} and \ref{ut:high} at the floor, its proved
physical coefficient bounds \eqref{ut:physical-coefficients}, and
the positive corner divisors above. At each physical step retain only
necessary inequalities for the current law and a complete closed
cover of its incoming domain. Every support or strip implication is
proved at the actual integer $n$, using its own incoming moment and
threshold conditions before its output is used.
Then every saved raw and defect output $(r_j,d_j)$ has the interpretation
\begin{equation}\label{ut:induction-invariant}
 \mathcal R_n\le r_j,\qquad nD\le Nd_j.
\end{equation}
The physical moment, endpoint, sign and threshold bounds exported
alongside it apply simultaneously to the same actual-$n$ law.
\end{proposition}
\begin{proof}
For the first pass there is no input $r_j$ or $d_j$.
Proposition~\ref{gc:initial-support} supplies the initial radii
$12/5$ on B0 and $13/5$ on B1 uniformly in the actual integer.
Before any defect feedback,
$M_4\le R\rho$ gives the respective ceilings $66/25$ and $78/25$;
$M_5\le RM_4$, $M_6\le RM_5$ and $|\gamma|\le\rho\le6/5$
license the triangle low bound and the universal or Cantelli high
bound. The initial signed-moment cover uses $\varphi\in[0,2/5]$,
the independent score interval and physical moment inequalities.
It assumes neither a central gap nor a final conditional moment cap.

The selected first raw cover maxima, with outward arithmetic, are
\[
 r_1^{\rm B0}=109743487/50000000,
 \qquad r_1^{\rm B1}=283144217/100000000.
\]
These are finite upper bounds on the raw prices in
\eqref{ut:raw-candidate}; their closed-domain arithmetic is part of
the companion calculation. Both are below $3<.4N\le nC\rho$.
It follows from \eqref{loc:conserved} that
\[
 \frac{\varphi}{6}(\gamma+3J)
            =C\rho+\sqrt n\mathcal L-\mathcal R_n/n>0.
\]
Consequently $\CE\rho-\phi(0)D/6>0$. Removing the Gaussian weight
in \eqref{loc:conserved} now lowers its left side and proves
\begin{equation}\label{ut:bootstrap}
 n(C-\CE)\rho+\kappa nD+n^{3/2}\mathcal L\le\mathcal R_n.
\end{equation}
Only at this point is the charged output from the same first cover
used. Its corner calculation, with $e_N=0$ and
\eqref{ut:excess-charge}, gives
$d_1=B_N/(\kappa N)$, with the saved outward values
\[
 d_1^{\rm B0}=129264343/10^8,\qquad
 d_1^{\rm B1}=16619657/10^7.
\]
Thus $nD\le Nd_1$, not merely $D\le d_1$.
Both $d_1<2< (\sqrt{10}-1)\rho_l$, so the next pass may use the
central gap. This order explains why calculating raw and charged
corner arrays in one routine is not a circular use of a defect cap.

Suppose \eqref{ut:induction-invariant} and the physical bounds have
been established at a later incoming stage. Since $n\ge N$, the
weaker physical statement $D\le d_j$ is available. All moment,
common-mean, gap and conditional-variance deductions can therefore
use $d_j$ in exactly their displayed physical directions.
For example, the additional p-local mean caps are
\[
 \mu^2\le p(1+\beta)/2,\quad
 \mu^2\le(1-p)(1-\beta)/2,\quad \mu^2\le p(1-p).
\]
They follow from signed-class Cauchy--Schwarz and retain the same
meaning on both probability halves. The endpoint-product distances
in \eqref{ut:gap-input} and all signed Pearson rectangles above
are likewise independent of $n$.

The physical inputs now license the next fixed Fourier candidates.
Their exact low and high expressions at the actual $n$ are bounded
by their floor-$N$ values, by the two lemmas. The floor density
interval \eqref{ut:incoming-envelope} contains the actual density;
the floor own-score interval \eqref{ut:own-score} contains the
actual score; and \eqref{ut:rebate} is a lower price for the actual
conserved phase. Equations~\eqref{ut:raw-candidate} and
\eqref{ut:charged-candidate} therefore hold pointwise in the actual
$H$. Their whole-cell corner maxima yield the next
$\mathcal R_n\le R_N$ and $nD\le B_N/\kappa$.
Writing the latter output as $d_{j+1}=B_N/(\kappa N)$ restores the
weighted invariant at every charged stage.

An excluded physical cell has an empty necessary interval or a
strictly negative upper bound for a nonnegative resource.
Any surviving law lies in at least one retained closed cell, so
maximizing its new bounds over all retained cells is valid.
Intersections with preceding bounds, finite minima of valid
candidates and outward rounding preserve the invariant.
This proves the assertion by induction over the prescribed finite
schedule, including repeated local INITIAL covers, the resource
subiterations and the final exports. The proof invokes neither
convergence of the iteration nor extremality at a smaller sample size.
\end{proof}

\subsection{Actual support and the interpretation of geometry}

The support hypotheses in Proposition~\ref{ut:transfer} use the
actual-index Gaussian majorant \eqref{loc:support-test}.
The original contact function $\mathcal D_n$ is zero on the full
support; the majorant need only be nonnegative there.
Write $s_n=\sqrt{n/(n-1)}$ and $e_n=n(s_n-1)-1/2$.
The possible extra allowance is
$n\rho s_n(C_{n-1}-C)_+=0$ under the assumed full predecessor.
A support ray or an excluded strip is obtained by proving the right
side strictly negative throughout that set, uniformly in all incoming
moment, score and threshold parameters. The Gaussian differences,
$s_n$ and $e_n$ in that calculation use the actual $n$.
For a bounded $H$ test, that bound on $H$ must already have been
proved by an earlier phase envelope. In particular the estimate
$e_n<.384/n$, proved in the retained derivation only for $n\ge36$,
is not used in these five actual-index requests.

The initial signed-moment interval follows from the support-free
inequality \eqref{ut:initial-signed-moment}; it requires no radius
or higher-moment feedback.

At the first support call the only moment restrictions are the
specified rho band and its consequences, together with the actual
score cap. An intermediate call consumes the moment intervals which
have already been proved by the preceding resource stage. Its new
radius licenses moment improvements only after the actual-index
strict contact test has passed. A final moment cap cannot license
the first support call, and a successful support calculation at $25$
cannot license a support or strip request at $26,\ldots,30$.
This is the one kind of acquisition step for which the floor-index
substitution is deliberately not made.

The distinction between weighted and physical quantities is also
necessary when exporting noise. From an acquired physical inequality
\[
 D\ge b+g\delta+h_NNs,\quad
 \delta=a_0+a_1,\quad b\ge0,\quad g,h_N>0,
\]
and $nD\le Nd_j$, put $r=N/n$ and $R_j=d_j-b$.
If $R_j>0$, $d_0=R_j/g$ and $L=R_j/h_N$ give
\begin{equation}\label{ut:export-plane}
 D\le rd_j,\qquad
 D\ge b+g\delta+rh_N\lambda_n,\qquad
 \frac{\delta}{rd_0}+\frac{\lambda_n}{L}
                 \le1-(n/N-1)b/R_j\le1,
 \quad\lambda_n=ns.
\end{equation}
This proves the units of the actual defect, inward intercept and noise
price. The joint intercept $L$ stays fixed. Without dividing by $R_j$,
the original inequality reads
$g\delta+rh_N\lambda_n\le rd_j-b$.
A negative right side excludes the cell, and a zero right side forces
$\delta=\lambda_n=0$. This also specifies the cases $R_j\le0$:
zero noise is retained only when the actual scaled right side is zero.

In the positive-residual case $R_j>0$, suppose physical support also gives $s\le B\delta$,
$B>0$, and $p(1-p)\ge v_l>0$. Retain first the full scaled
plane in \eqref{ut:export-plane} and put
\[
 c=1-(n/N-1)b/R_j,\qquad
 H_N=\left(\frac1L+\frac1{NBd_0}\right)^{-1}.
\]
The same scalar elimination as in \eqref{loc:harmonic-consequence},
now using $\delta\ge\lambda_n/(nB)$, $h=rh_N$,
$\overline D=rd_j$ and $nr=N$, gives
\begin{equation}\label{ut:scaled-harmonic}
 \lambda_n\left(\frac1L+\frac1{NBd_0}\right)\le c,
 \qquad \lambda_n\le cH_N.
\end{equation}
This is the stronger scalar export used by the actual-index adapter.
For a nonempty cell $c\ge0$, and $b\ge0$ gives $c\le1$;
with this same support constant $B$, $H_N$ is therefore also
a conservative scalar cap. At $c=0$ the scaled plane retains
$\delta=\lambda_n=0$.

If only the deliberately weakened, unscaled joint plane
$\delta/d_0+\lambda_n/L\le1$ is retained, it gives
\[
 \lambda_n\le\left(\frac1L+\frac1{nBd_0}\right)^{-1},\qquad
 \theta\le\left(1+\frac{nv_l}{L}
                         +\frac{v_l}{Bd_0}\right)^{-1}
 \le\left(1+\frac{Nv_l}{L}
                         +\frac{v_l}{Bd_0}\right)^{-1}.
\]
These weaker premises license the displayed actual-$n$
scalar expression and the transported variance fraction. They alone
do not license replacing that scalar expression by $H_N$.
The stronger cap in \eqref{ut:scaled-harmonic} is available when the
full scaled plane is retained. Both harmonic expressions are separate
scalar consequences; the original simultaneous intercept in
\eqref{ut:export-plane} remains $L$.
All further signed refinements may use the transported physical
$\theta$ and the actual support, followed by a new application of
the same joint-defect inequality. This completes the mathematical
transfer needed for the five actual-index acquisition schedules.

% Formula sources read for this appendix (all frozen source files are unchanged):
% source/analytic/UNIFORM25_CANTELLI_PRODUCTGAP_REVIEW.txt;
% source/analytic/UNIFORM25_BETA0_SOURCE_AND_COMPOSITION_REVIEW.txt;
% source/cutoff25_20261011/new_root/initial25_b1_exact_v1/
% RETAINED_ANALYTIC_FRAMEWORK.tex, initial34.py, kernel.py, high_window_core.py;
% the B0 initial25_b0_v2 counterparts; resource25_b0_v1 and resource25_b1_v1/
% candidate_charged_score_v2.py, contracted_kernel.py, general_paired_exact.py,
% corrected_phase_callback_extended.py, resource25_core.py.

\section{Gaussian contact as a Peano remainder}\label{gc:section}

The initial support bounds admit a single analytic proof. The key is to
recognize the contact kernel as a Gaussian Hermite remainder, and then
apply the classical Peano-kernel representation. This replaces the
initial actual-index support calculations and supplies the radial
contractions used by the tail route. Remaining actual-index refinements
and coordinate-strip exclusions retain their own proofs.

Let $Z$ be standard normal. For $g\in C^3(\mathbb R)$ with bounded
third derivative, put
\[
 L_xg=g(x)-\E g(Z)-x\E g'(Z)-\frac{x^2-1}{2}\E g''(Z).
\]
These expectations exist because $g,g',g''$ grow at most cubically,
quadratically and linearly.

\begin{lemma}[A uniform Gaussian Peano kernel]\label{gc:kernel}
For $|x|\ge2$,
\begin{equation}\label{gc:kernel-bound}
 |L_xg|\le\|g'''\|_\infty
       \left\{\frac{|x|^3-3|x|}{6}+\frac34\right\}.
\end{equation}
\end{lemma}
\begin{proof}
Taylor's integral remainder anchored at zero, followed by Fubini, gives
\begin{equation}\label{gc:peano}
 \begin{aligned}
 L_xg&=\int_{\mathbb R}g'''(t)K_x(t)\,dt,\\
 K_x(t)&=\frac{(x-t)^2}{2}
       \{\mathbf1_{\{t<x\}}-\overline\Phi(t)\}+(t/2-x)\phi(t).
 \end{aligned}
\end{equation}
Indeed $L_x$ annihilates all quadratic polynomials. Applying it to
$(u-t)_+^2/2$, and using the first two truncated Gaussian moments,
gives exactly this kernel. The anchored remainders are bounded by
Gaussian absolute moments of orders at most three, justifying the
interchanges. The displayed kernel has integrable Gaussian tails.
Testing $g(u)=u^3/6$ and reflecting give
\[
 \int K_x(t)\,dt=\frac{x^3-3x}{6},\qquad
 K_{-x}(-t)=-K_x(t)\quad\text{almost everywhere}.
\]

Take $x\ge2$. We first show $K_x(t)\ge0$ for $t\le x/2$.
For $t=-y\le0$, this is the Mills inequality
\[
 m(y):=\frac{\overline\Phi(y)}{\phi(y)}
 \ge\frac{2x+y}{(x+y)^2}.
\]
Cauchy--Schwarz for the integrals of $1,u,u^2$ against
$e^{-yu-u^2/2}du$ gives $m^2+ym\ge1$. Hence
\[
 m(y)\ge\frac{\sqrt{y^2+4}-y}{2}
       \ge\frac{y+4}{(y+2)^2}
       \ge\frac{y+2x}{(y+x)^2}.
\]
The middle comparison follows by squaring positive quantities; the
last follows by differentiation in $x\ge2$.

For $0\le t\le1$, $K_x(t)$ increases in $x\ge2$, since
$\partial_xK_x=(x-t)\Phi(t)-\phi(t)>0$. The elementary bounds
$39/100<\phi(0)<2/5$, $e^{-u}\ge1-u$ and
$e^{-u}\le1-u+u^2/2$ give $K_2(t)\ge P(t)$, where
\[
 P(t)=\frac15-\frac t{50}-\frac{13t^2}{100}
       -\frac{7t^3}{200}+\frac{3t^4}{100}-\frac{3t^5}{400}.
\]
On $[0,1]$, $P'(t)\le-49t/200-1/50<0$, so
$P(t)\ge P(1)=3/80>0$.
For $1\le t\le x/2$ the same monotonicity in $x$ gives
\[
 K_x(t)\ge K_{2t}(t)
       =\frac t2\{t\Phi(t)-3\phi(t)\}>0.
\]
The expression in braces increases for $t\ge1$; at $1$ it is
positive because $\Phi(1)\ge33/40$ and $\phi(1)<49/200$.

Writing $a=x/2$, we can now bound the negative part by
\begin{align}\label{gc:negative-part}
 \int K_x^-&\le\int_a^{2x}(x-t/2)\phi(t)\,dt
       +\frac12\int_x^\infty(t-x)^2\overline\Phi(t)\,dt\notag\\
 &\le\frac32\phi(a)+\frac12\phi(2x)+\frac{\phi(x)}{x^4}
 <\frac{5943}{16000}<\frac38.
\end{align}
The first integral uses $\overline\Phi(a)\le\phi(a)/a$.
The second is $\E(Z-x)_+^3/6$, at most $\phi(x)/x^4$ by
$\phi(x+u)\le\phi(x)e^{-xu}$. The last rational bound uses
$\phi(1)<49/200$, $\phi(2)<11/200$ and $\phi(4)<1/1000$.
Each follows from a positive partial sum of the exponential series.
Thus $\|K_x\|_1=\int K_x+2\int K_x^-$ proves
\eqref{gc:kernel-bound}.

\end{proof}

The signed moment needed for contact also has a support-free bound.
For $a=\E|X|>0$, $\beta=\E[X|X|]$ and $\rho=\E|X|^3$,
\begin{equation}\label{ut:initial-signed-moment}
 a\rho\ge1+\beta^2,\qquad a^2+\beta^2\le1,\qquad
 |\beta|\le\sqrt{\frac{\rho\sqrt{\rho^2+8}-\rho^2-2}{2}}.
\end{equation}
For the first inequality, apply Bessel's inequality to $|X|^{3/2}$
and the orthogonal functions $|X|^{1/2}$ and
$\operatorname{sgn}(X)|X|^{1/2}$. Their squared norms are $a$ and
their inner products with $|X|^{3/2}$ are $1,\beta$.
Cauchy--Schwarz for $X$ and $|X|-a$ proves the second inequality;
eliminating $a$ proves the last.

\begin{proposition}[Initial support without index subdivisions]\label{gc:initial-support}
Let $(P,H)$ be a standardized positive closed global maximizing pair
at index $n\ge25$, with $C=C_n\ge\CE$ and $C_{n-1}\le C$.
Then
\[
 \rho\le11/10\ \Longrightarrow\ \operatorname{supp}P\subset(-12/5,12/5),
\]
\[
 \rho\le6/5\ \Longrightarrow\ \operatorname{supp}P\subset(-13/5,13/5).
\]
No higher moment or defect bound is required.
\end{proposition}
\begin{proof}
For $g(u)=n^{3/2}\Phi((H-u)/\sqrt{n-1})$, Gaussian convolution gives
\begin{align*}
 \E g(Z)&=n^{3/2}\Phi(H/\sqrt n),\\
 \E g'(Z)&=-n\phi(H/\sqrt n),\qquad
 \E g''(Z)=-H\phi(H/\sqrt n).
\end{align*}
Thus $L_xg=\mathcal H_n(H,x)$ exactly. Also
$\|g'''\|_\infty=(n/(n-1))^{3/2}\phi(0)<213/500$ for $n\ge25$.
For $s_n=\sqrt{n/(n-1)}$ and $e_n=n(s_n-1)-1/2$,
\begin{equation}\label{gc:scale}
 0<e_n<\frac{2}{5n}\le\frac2{125}.
\end{equation}
Indeed, for $z=1/n\le1/25$,
\[
 (1-z)(1+z/2+2z^2/5)^2-1
 =z^2\{1/20-13z/20-6z^2/25-4z^3/25\}>0.
\]
The bracket decreases and is positive at $1/25$; omitting the
quadratic correction proves the lower bound.

By \eqref{ut:initial-signed-moment}, $\rho\le11/10$ gives
$|\beta|<127/500$, and $\rho\le6/5$ gives $|\beta|<44/125$.
For either $(\rho_u,b,R)=(11/10,127/500,12/5)$ or
$(6/5,44/125,13/5)$, put
\begin{align*}
 Q(r)&=-r^3+3\rho_ur^2/2+3br+2\rho_u/125,\\
 V(r)&=\frac{213}{500}\{(r^3-3r)/6+3/4\}
                                      +\frac{819}{2000}Q(r).
\end{align*}
At $r=R$, $Q<0$, $Q'<-8$ and $Q''<-11$; since $Q'''=-6$,
the multiplier of $C$ remains negative on the whole ray.
The contact inequality \eqref{loc:support-test} therefore has
right side at most $V(r)$ for $|x|=r\ge R$, using
$C\ge\CE>819/2000$ in the negative term. Direct substitution gives
\[
 V(12/5)=-278919/1250000,\qquad V(13/5)=-85993/1250000.
\]
In both rows $V'(R)<-5/2$, $V''(R)<-7/2$ and $V'''<-2$.
Hence $V<0$ on both entire rays, contradicting contact at any
support point there, including atoms and accumulation points.
\end{proof}

\begin{corollary}[All-index radial support at zero predecessor loss]\label{gc:tail-support}
Let $(P,H)$ be a standardized positive closed global maximizing pair
with $n\ge47$, $C=C_n\ge\CE$ and $C_{n-1}\le C$.
For each row below, the already established moment bounds
$\rho\le\rho_u$, $|\beta|\le b$ imply
$\operatorname{supp}P\subset(-R,R)$:
\[
\begin{array}{c|c|c|c}
 \rho_u&b&R&\text{contact upper bound at }R\\\hline
 7361/4000&643/1000&96/25&-1/20\\
 139/100&43/100&3&-23/100\\
 12563/10000&1963/5000&27/10&-9/500\\
 603687/500000&4014/15625&5/2&-49/500
\end{array}
\]
\end{corollary}
\begin{proof}
The predecessor allowance is zero. The same Gaussian function has
$\|g'''\|_\infty<207/500$ and $e_n<2/235$ for $n\ge47$.
For each row put
\begin{align*}
 Q(r)&=-r^3+3\rho_ur^2/2+3br+2\rho_u/235,\\
 V(r)&=\frac{207}{500}\{(r^3-3r)/6+3/4\}
       +\frac{819}{2000}Q(r).
\end{align*}
Direct substitution gives the strict upper bounds in the table.
In all four rows $Q(R)<0$, $Q'(R)<-8$, $Q''(R)<-11$,
$V'(R)<-5/2$, $V''(R)<-7/2$ and $V'''<-2$.
The whole-ray argument in Proposition~\ref{gc:initial-support}
therefore applies unchanged.
\end{proof}

These moment bounds are used in predecessor order. The global envelope
supplies the first row. The original central preparation, performed
at radius $3.84$, supplies the last row for all nine surviving central
intervals. The complementary acquisition at $3.84$ supplies the second
row; after contraction to $3$, its left-input and signed-resource
calculations supply the third. Thus no contraction uses a moment
bound obtained from its own output radius. The later coordinate-strip
exclusions retain their original bounded-threshold proofs.
The unsigned $\beta$ caps in these resource rows follow from
$\beta=(1-\theta)(1-2p)+\theta_+-\theta_-$, hence
$|\beta|\le(1-\theta)|1-2p|+\theta$; no sign of $\beta$ is assumed.

\section{The analytic profile parameters}\label{app:parameters}
The table gives the original joint inward intercept $\Delta$ and
noise intercept $\Lambda$, together with the improved necessary
actual-noise cap $L_*$, for the $34$ middle-moment profiles.
All entries are exact terminating decimals. The simultaneous
signed coefficients, outward support caps and contact domains
are supplied as exact rational numbers in the accompanying
\texttt{profile-parameters.json}; they are acquired conditions
on an extremizer, not independent free parameters.
\begingroup\small
\begin{longtable}{ccrrr}\toprule
$n$ & interval & $\Delta$ & $\Lambda$ & $L_*$\\\midrule\endhead
$26$ & $I_{34}$ & 0.07393 & 0.448664 & 0.26204965\\
$26$ & $I_{35}$ & 0.078697 & 0.469161 & 0.27350801\\
$26$ & $I_{36}$ & 0.08286 & 0.484491 & 0.28185367\\
$26$ & $I_{37}$ & 0.086366 & 0.494588 & 0.29863673\\
$26$ & $I_{38}$ & 0.083933 & 0.467679 & 0.28297254\\
$26$ & $I_{39}$ & 0.077665 & 0.420047 & 0.2556063\\
\addlinespace
$27$ & $I_{34}$ & 0.071192 & 0.448664 & 0.25573091\\
$27$ & $I_{35}$ & 0.075782 & 0.469161 & 0.26738858\\
$27$ & $I_{36}$ & 0.079791 & 0.484491 & 0.2759327\\
$27$ & $I_{37}$ & 0.083167 & 0.494588 & 0.29258472\\
$27$ & $I_{38}$ & 0.080824 & 0.467679 & 0.27794155\\
$27$ & $I_{39}$ & 0.074788 & 0.420047 & 0.25182702\\
$27$ & $I_{40}$ & 0.06347 & 0.34526 & 0.20897126\\
\addlinespace
$28$ & $I_{34}$ & 0.068649 & 0.448664 & 0.25004307\\
$28$ & $I_{35}$ & 0.073076 & 0.469161 & 0.26192751\\
$28$ & $I_{36}$ & 0.076941 & 0.484491 & 0.27068902\\
$28$ & $I_{37}$ & 0.080197 & 0.494588 & 0.28726658\\
$28$ & $I_{38}$ & 0.077938 & 0.467679 & 0.27351126\\
$28$ & $I_{39}$ & 0.072117 & 0.420047 & 0.24847935\\
$28$ & $I_{40}$ & 0.061204 & 0.34526 & 0.20687478\\
\addlinespace
$29$ & $I_{34}$ & 0.066282 & 0.448664 & 0.24487077\\
$29$ & $I_{35}$ & 0.070556 & 0.469161 & 0.25700268\\
$29$ & $I_{36}$ & 0.074288 & 0.484491 & 0.26599549\\
$29$ & $I_{37}$ & 0.077432 & 0.494588 & 0.28254209\\
$29$ & $I_{38}$ & 0.07525 & 0.467679 & 0.26957152\\
$29$ & $I_{39}$ & 0.06963 & 0.420047 & 0.24548928\\
$29$ & $I_{40}$ & 0.059093 & 0.34526 & 0.20498657\\
\addlinespace
$30$ & $I_{34}$ & 0.064073 & 0.448664 & 0.24012598\\
$30$ & $I_{35}$ & 0.068204 & 0.469161 & 0.25252129\\
$30$ & $I_{36}$ & 0.071812 & 0.484491 & 0.26175592\\
$30$ & $I_{37}$ & 0.074851 & 0.494588 & 0.27830578\\
$30$ & $I_{38}$ & 0.072742 & 0.467679 & 0.26603838\\
$30$ & $I_{39}$ & 0.067309 & 0.420047 & 0.24279925\\
$30$ & $I_{40}$ & 0.057124 & 0.34526 & 0.20327613\\
\bottomrule\end{longtable}\endgroup
The analytic programs use the Python 3 standard library and retain
exact rational arithmetic. In the companion directory
\path{finite-steps/source/analytic}, the seven actual-$30$
profiles can be recomputed with
\par\noindent\texttt{python -B}\ \path{replay_n30_analytic_panels_20261027.py}
\ \texttt{--output fresh.json}\par
An explicit \texttt{--output} argument selects a fresh filename;
the program refuses to overwrite an existing result.
The option \texttt{--indices 37} recomputes just that complete
profile, including its contact domain and the CDF bridge.
The finite-index mathematical companion states the count and phase
composition for every profile; the executable calculation supplies
the enclosed faces, derivatives and strict margins used there.

\section{The seven lower-moment comparison plans}\label{app:finite-details}

This appendix gives the complete plans underlying Proposition~\ref{fc:lowprofiles}.
All use the common comparison, curvature and allocation rules of the main text.

\begin{lemma}[One lower-moment profile]\label{fc:low40}
Under the retained actual-$26$ acquisition and the global predecessor
$C_{25}\le\CE$, the entire original profile
$p\in[.40,.41]$, $1\le\rho\le11/10$ has a strictly negative
conditional count residual. A uniform upper bound is
$-5610547/10^{12}$.
\end{lemma}
\begin{proof}
The original inward and noise intercepts are
$\Delta=.057809$ and $\Lambda=.301161$.
The necessary signed resource plane improves the actual-noise cap
to $L_*=.17805629$. The original contact inequalities exclude
$25$ of the $27$ count roots, leaving only $10$ and $11$.

Current-count comparisons cover count $10$ on $[-1/2,0]$ with
slope multiplier $1$, and count $11$ on its two half intervals
with multipliers $3/2$ and $1$. Their upper bounds are respectively
$-.000005610547$, $-.000934723659$, and $-.003337615615$.

For the positive half of count $10$, retain the seven noise
intervals and four allocation intervals above.
On the first noise interval join the current price to the
reflected-square price at $t=1/32$; on the second join at
$1/8$; on intervals $3$--$7$ join at $1/4$.
Each reflected-square price uses reference count $11$, width
$h=7/8$ and Corollary~\ref{ce:allocation-interval}.
Here $11>26p$, $p<1/2$, and the positive affine inward cap
decreases in allocation. The moment floor therefore increases
in allocation. All next-atom, legal-width and same-noise derivative
gates hold on their whole stated domains. In the first interval
$y_{\max}=3/32>0$ is allowed by the extended variance-derivative
bound in Section~\ref{pp:section}.

Every current prefix is at most $-.000025632850$;
every clipped suffix is at most $-.001597137326$.
Probability and phase chords include every closed boundary.
The two pieces in each noise/allocation region meet at the stated
phase, and the full noise union is $[0,L_*]$.
Together with the three half-interval comparisons and the contact
exclusions this covers the whole original profile. The least strict
bound is $-.000005610547$.
\end{proof}

This replaces a conditional count comparison. The original support,
floor and acquisition proofs remain its inputs; replaying the pure
physical acquisition formulas does not establish those inputs anew.

\begin{lemma}[An adjacent lower-moment profile]\label{fc:low39}
Under the retained actual-$26$ acquisition and $C_{25}\le\CE$,
the entire original profile $p\in[.39,.40]$, $1\le\rho\le11/10$
has conditional count residual at most $-98710877/10^{12}$.
\end{lemma}
\begin{proof}
The original intercepts are $\Delta=.05769$, $\Lambda=.30708$;
the same signed resource argument gives $\lambda\le .17375152$.
Contact excludes $24$ count roots. Only $9,10,11$ remain, and
the count-$9$ contact lower phase is greater than $1/4$.
Count $11$ has the two current-count half-interval comparisons
with slope multipliers $3/2,1$.

For the lower counts one reference law covers both sides of a
nearest-count boundary. Use source count $9$, reference count $10$
and width $7/8$. On noise intervals $1,2$ its auxiliary phase runs
through $[1/4,7/8]$, and on intervals $3$--$7$ through $[1/4,3/4]$.
The exact threshold identity
\[
 9+(1+t)=10+t
\]
therefore covers count $9$ and the lower part of count $10$ by
one continuous price in each noise group. Lemma~\ref{mc:staircase}
holds for every real phase, so an auxiliary phase greater than
$1/2$ is permitted. Here $10<26p$ throughout the probability
interval, and Corollary~\ref{ce:allocation-interval} selects the
lower allocation endpoint for both variance and moment cost.

The remaining count-$10$ negative phases are $[-1/8,0]$ on
the first two noise intervals and $[-1/4,0]$ on the others;
current-count prices cover them. Its positive half uses the
same $1/32,1/8,1/4$ handoffs as in Lemma~\ref{fc:low40}, now
with this profile's own resources. Reference count $11>26p$
selects the upper allocation endpoint for variance and the lower
one for moment cost. All legal-width, next-atom, noise-derivative
and continuous probability/phase gates are strict on their stated
domains. The whole closed unions retain every original count,
noise and allocation region, with largest upper bound
$-.000098710877$.
\end{proof}

\begin{lemma}[Reversed signed noise prices]\label{fc:low41}
Under the retained actual-$26$ acquisition and $C_{25}\le\CE$,
the entire original profile $p\in[.41,.42]$, $1\le\rho\le11/10$
has conditional count residual at most $-24705293/(5\cdot10^{11})$.
\end{lemma}
\begin{proof}
The original intercepts are $\Delta=.054236$, $\Lambda=.276951$.
Here $h_+-h_-<0$. The positive endpoint coefficients in
\eqref{loc:two-endpoint-noise} still give a valid finite resource
iteration and $\lambda\le L_*=.162619$, retaining the original
joint intercept $\Lambda$. For the actual moment floor,
\eqref{ce:signed-allocation-gate} is greater than $.473372164977$;
both negative terms from the increasing inward cap are included.
Thus $G$ increases with allocation on the whole domain.

For current-count prices, the necessary weaker plane
\[
 g(a_0+a_1)+\lambda h_-
 \le\mathfrak R+L_*(h_--h_+)_+
\]
supplies an allocation-independent inward cap. Its intercept is
smaller than the original $\Delta$, and its noise slope is larger
than $\Delta/\Lambda$, as required by the fixed-reference estimate.

Contact excludes $24$ count roots and leaves $10,11,12$.
Count $10$ starts above $-1/8$, and count $12$ ends below $-1/4$.
Current-count prefixes cover count $10$ through $1/32,1/8,1/4$
on the first, second and remaining noise intervals.
One source-$10$, reference-$11$ square then covers both its
positive suffix and the negative prefix of count $11$.
On noise intervals $1,2$ the width is $7/8$ and its auxiliary
phase ends at $7/8$. On intervals $3$--$6$ use width $9/8$
with phase $[1/4,7/8]$. On the last interval use width $7/8$
and phase $[1/4,3/4]$. The identity $10+(1+t)=11+t$
gives the original count-$11$ phases exactly.

Current-count prices complete count $11$ through zero and on its
positive half; count $12$ uses slope multiplier $3/2$ on
$[-1/2,-1/4]$. The original seven noise intervals and four
allocation intervals are unchanged. All width, next-atom,
allocation and noise gates, followed by the probability/phase
chords, give a strict bound on the full closed unions.
Their largest upper bound is $-.000049410586$.
\end{proof}

\begin{corollary}[A complete lower-moment strip]\label{fc:p41-strip}\label{fc:n27-low41}\label{fc:n28-low41}\label{fc:n29-low41}\label{fc:n30-low41}
At each integer $n=26,\ldots,30$, under the retained actual-$n$ acquisition
and $C_{n-1}\le\CE$, the original profile $p\in[.41,.42]$,
$1\le\rho\le11/10$ has the following conditional count bounds.
The original noise intercept is $\Lambda=.276951$ in all rows.
The five residual bounds are the corresponding rows of
Proposition~\ref{fc:lowprofiles}.

\end{corollary}
\begin{proof}
Use the same comparison rules with each integer's own masses,
variance and moment normalization. Lemma~\ref{fc:low41} gives $26$.
At $27$, contact excludes
$25$ of the $28$ roots and leaves $10,11,12$.
Reference count $11<27p$ uses lower allocation
endpoints for variance and cost, while $12>27p$ uses the two
allocation corners. The signed moment derivative gate is positive.
On the lowest noise interval the source-$10$ square
runs through auxiliary phase $15/16$, which is exactly count-$11$
phase $-1/16$. The current price on $[-1/16,1/32]$ and its
clipped-square continuation then close the middle count.

At $28$, the remaining counts are $11$ and $12$.
Their outer phase intervals are $[-17/64,1/2]$ and $[-1/2,15/32]$.
The same fixed current
references for count $11$ use the phase faces $-1/8$ and $0$,
followed by the square continuations. Count $12$ joins current
slopes $3/2$ and $1$ at phase $-1/32$. Reference count $12>28p$
has fixed variance direction throughout the probability interval.

At $29$, contact excludes $27$ roots and leaves $11,12,13$.
Reference count $12$ crosses $29p$ at $p=12/29$ inside the
original probability interval. Use both variance-endpoint branches
in \eqref{ce:variance-endpoint-max}, with the same actual noise
and the lower-allocation moment cost in both. Each branch has
its own whole-interval noise and curvature checks before their
maximum is bounded. The lowest-count square starts at phase
$1/8$, below the contact lower bound $.15468890249$; on the
sixth noise interval, auxiliary phase $13/16$ meets original
count-$12$ phase $-3/16$.

At $30$, contact leaves only $12,13$, with valid outer phase
intervals $[-7/16,1/2]$ and $[-1/2,5/16]$. The negative part
of count $12$ uses a lower-allocation square and current prices;
on the lowest noise interval the same current reference is
interpolated through one extra phase face, $-1/8$.
Its positive part hands off at $1/128$ with width $2/3$ on
that interval. Here $y_{\max}=125/384>0$ and the one-clip
and variance-derivative conditions hold on the whole region.
Higher-noise square prices use width $5/4$ where specified
by the plan, and the shifted-count joins close count $13$.

All these explicit finite plans retain the seven original noise and
four allocation intervals. Their remaining pieces use the current
or square envelopes, and their closed unions cover every retained
contact interval. Their largest upper bounds are those in the table.
\end{proof}

\section{Locations of the companion proofs}\label{app:companions}

The moment-defect and original count arguments are due to
Rodziewicz--Ko\l{}odziejek \cite{RK2026}; the completion notes identify
their reconstruction scope. The companion sources retain their original
attribution, theorem labels and statement scopes. A label quoted below belongs to the named companion,
not to the numbering of this article. The source collection is prepared
by \path{prepare_global_companions.py}; it includes the local TeX
inputs of the multifile large-index proof.

\paragraph{Input (I).}
In \path{companions/low-index/}, the first-index bound is
\path{n1/n1_base.tex}, Theorem~1 (\texttt{thm:base}).
The source \path{n2-n3/manuscript.tex} contains
\texttt{n2:thm:main} and \texttt{n3:thm:main}, together with their
support reductions and continuous-domain polynomial certificate
implications. The complete fourth-index proof is available in the
existing public manuscript PDF
\path{companions/references/rk-v3/moment-defects-v3.pdf},
Section~26, Theorem~26.1 and its proof. A TeX source for that complete
fourth-index proof is not included.

The subsequent compositions are Theorem~1
(\texttt{composition}) in \path{n5/five_sum_completion_20261025.tex};
Proposition~1 and its completed composition in each of
\path{n6/six_sum_interface_20261025.tex},
\path{n7/seven_sum_interface_20261025.tex},
\path{n8/eight_sum_interface_20261025.tex}, and
\path{n9/nine_sum_interface_20261025.tex};
the theorem entitled \emph{Unrestricted tenth-index conclusion}
in \path{n10/main.tex}; and the ordered implication at the end of
\path{n11-14/conditional_11_14_interface_20261026.tex}
(which has no separately numbered conclusion).
The final blocks are Theorem~1 in \path{n15-19/main.tex},
Proposition~1 and its completed composition in
\path{n20/twenty_sum_interface_20261025.tex},
Theorem~1 in \path{n21-24/main.tex}, and
Theorem~1 in \path{n25/main.tex}.
The high-moment complement at the actual indices $25$--$30$
is stated and derived in \path{n25-fourier/main.tex}.

These interfaces are accompanied by continuous derivations.
The source
\path{companions/common/finite-analytic/BE_cutoff35_streamlined.tex}
contains \path{global:attainment}, \texttt{global:contact},
\path{complete44ent:smooth-fourier},
\path{complete44ent:fourier-early-low},
\path{complete44ent:re-enclosure}, \texttt{res:contract},
and the section \emph{One finite-count comparison for all the new indices}.
The additional small-index proofs retain their original text format under
\path{companions/low-index/continuous-sources/work/}.
In particular, the fifth-index
\path{cutoff5_20261020/audit/ANALYTIC_DEPENDENCIES_V2.txt}
locates the conserved-identity, finite-window LOW, resource and count
derivations; their mathematical proof texts are included alongside that
source map. They are distinct from numerical success records.
The original frozen arithmetic inputs and actual returned objects are
retained under \path{companions/low-index-support/}.
Its \path{read_low_index_support.py} describes the existing joins,
their exact input locations and their Python or compiler requirements.

\paragraph{Input (II).}
In \path{companions/actual-index/}, the coarse entry from
$C_{30}<\CE+.01$ and the literal bridges to $34$ are
Propositions \texttt{prop:31} and \texttt{prop:bridges} in
\path{n31-34/main.tex}. The exceptional band-$7$ defect-domain
argument is \path{n31-domain/sharp31_B7_D_domain_20261025.tex}.
Continue with the theorem entitled
\emph{Unrestricted index--35 conclusion} in
\path{n35/actual35_conditional_composition_20261025.tex},
the proposition entitled \emph{Unrestricted conclusions at 36, 37 and 38}
in \path{n36-38/actual36_38_conditional_composition_20261026.tex},
and Propositions~1 and~2 in \path{n39/main.tex}.
The actual-$39$ continuous-domain extension is separately included in
\path{n39-continuous/actual39_continuous_lemmas_20261026.tex}.
Theorem \texttt{thm:main} in \path{n40-48/main.tex} contains the
bridge through $48$; the present composition uses its steps only
through $46$. Its two local implications are
\texttt{prop:forty} and \texttt{prop:two}.
The directories \path{n40-continuous/}, \path{n40-original/},
\path{n42-continuous/}, and \path{n42-local/} contain the corresponding
actual-index and finite-band continuous arguments.
The corresponding compressed inputs, returned objects and existing
join programs are retained under \path{companions/bridge-support/}.
The accompanying \path{read_bridge_support.py} locates the mathematical
objects and launches the original joins against those collected inputs.

\paragraph{Input (III).}
The local resource implication \path{re:localG} in
\path{paper_entry/resource_entry.tex} is used with $b\ge0$, $P\ge0$,
$\ell_-\ge0$ and $L_0\ge0$. These are the physical coefficients of
\path{conf:moment} and the nonnegative noise subintervals; they license
the substitutions $nP\ge NP$ and $n\theta\ge\ell_-^2L_0$.

The proposition \texttt{global:local} in
\path{companions/large-index/analytic-route/paper/sections/global_route.tex}
contains the zero-allowance uniform local exclusion (III) used here.
In the present proof, Corollary~\ref{gc:tail-support} replaces its
initial radius-$3.84$ calculation over six index bands, both left radial
contractions and the central radius-$2.5$ contraction. The resource
schedule, coordinate-strip exclusions and local comparisons are
retained. The broader original heat-expansion route remains in the archive.
The continuous analytic route is the multifile source
\path{companions/large-index/analytic-route/paper/main.tex}:
\texttt{re:acquired}, \texttt{re:enclosure}, and
\path{shortloc:theorem} supply its resource acquisition and
local comparison. Its included coefficient atlas derives
compensated inversion and the conserved identity, rather than treating
them as numerical premises. The independently completed $n\ge49$
route is also retained in
\path{companions/large-index/completion/n49_completion_20261025.tex};
the new induction from $45,46$ uses the local exclusion directly.
The selected compressed source/input archive and its original returned
calculation archive are included under \path{companions/uniform-tail/}.
The accompanying \path{prepare_uniform_tail_support.py} can unpack
them to a new directory, reconstruct the exact local cover, and run
the original transcript-to-target join. Its Windows schedule launcher
changes two path-key expressions to POSIX spelling; all mathematical
conditions and original archive contents are preserved.

\paragraph{Weighted localization and the 34 finite profiles.}
Appendix~\ref{ut:section} gives the mathematical proof of
\eqref{loc:weighted-transfer}, including the raw first step, all LOW and
HIGH terms, the weighted induction and the order of actual-index support
acquisition. The earlier fixed-floor derivations are retained in
\path{companions/localization/analytic/UNIFORM25_CANTELLI_PRODUCTGAP_REVIEW.txt}
and
\path{companions/localization/analytic/UNIFORM25_BETA0_SOURCE_AND_COMPOSITION_REVIEW.txt}.
The inward-plane scaling is derived in
\path{companions/localization/original/new_audit/SCALED_INWARD_PHYSICAL_DERIVATION_V1.txt}.
Their continuous foundation is the resource-enclosure and simultaneous-contract
argument identified above; actual-index support acquisitions are still
separate inputs.

The initial B0/B1 radii now follow directly from
Proposition~\ref{gc:initial-support}. Their original actual-index
radius computations are no longer premises of this first step;
later support refinements still use their own incoming data.
The Peano-kernel and rational-ray calculations are also retained in
\path{gaussian_contact_peano_20261028.tex} and
\path{gaussian_contact_peano_20261028.py} in the finite analytic directory.

The full finite comparison text is
\path{companions/finite-steps/source/analytic/finite_gap_shared_stoploss_20261026.tex}.
Theorems \path{global-c26-complete}, \path{global-c27-complete},
\path{global-c28-complete}, \path{global-c29-complete}, and
\path{global-c30-complete} give the five actual-index compositions.
Their proofs identify the constituent profile lemmas.
At $26$, these include \path{p34-p35-count8-complete},
\path{p35-count9-second-moment-complete},
\path{p36-p37-count10-second-moment-complete}, and the complete-panel
corollaries for $36$--$39$.
At $27$, the profile conclusions are
\path{n27-p34-complete}, \path{n27-p35-complete},
\path{n27-p36-complete}, \path{n27-central-complete}, and
\path{n27-p39-complete}, \path{n27-p40-complete};
the central corollary covers both $37$ and $38$.
The corresponding $28$ statements are
\path{n28-p34-p35-complete}, \path{n28-p36-complete},
\path{n28-p37-complete}, \path{n28-p38-complete}, and
\path{n28-p39-p40-complete}.
The last two groups are the lemmas
\path{n29-seven-panels-complete} and
\path{n30-seven-panels-complete}.
The lower-moment replacement in Lemma~\ref{fc:low40} is the
separate program \path{probe_b0_shared_reference_20261027.py}
in that analytic source directory, with the exact output
\path{result_b0_n26_p40_complete_frozen_20261027.json}.
Its \texttt{--complete} entry recomputes this whole conditional
profile without reading a stored numerical result.
Lemma~\ref{fc:low39} uses the separate
\path{probe_b0_p39_shared_20261027.py} and
\path{b0_allocation_interval_clip_20261027.py}, with output
\path{result_b0_n26_p39_complete_frozen_20261027.json} in the
same analytic directory.
The signed-price extension in Lemma~\ref{fc:low41} is
\path{probe_b0_p41_shared_20261027.py}, with
\path{b0_two_endpoint_budget_20261027.py} and
\path{b0_signed_allocation_interval_20261027.py}; its result is
\path{result_b0_n26_p41_complete_frozen_20261027.json}.
Corollary~\ref{fc:n27-low41} reuses these helpers through the
explicit \texttt{PLAN} in
\path{probe_b0_n27_p41_shared_20261027.py}; the actual-$27$ output is
\path{result_b0_n27_p41_complete_fresh_20261027.json}.
The actual-$28$ plan is
\path{probe_b0_n28_p41_shared_20261028.py}, with output
\path{result_b0_n28_p41_complete_fresh_20261028.json}.
The crossing application at $29$ is the \texttt{PLAN} in
\path{probe_b0_n29_p41_shared_20261028.py}, using
\path{b0_variance_endpoint_branches_20261028.py}; its output is
\path{result_b0_n29_p41_complete_fresh_20261028.json}.
The actual-$30$ plan is
\path{probe_b0_n30_p41_shared_20261028.py}, with final output
\path{result_b0_n30_p41_complete_frozen_20261028.json}.
Its earlier failed complete candidate is retained in the research
sources and supplies no completed profile bound.
The finite lower-envelope implementation in
\path{probe_classical_order_compression_20261028.py} replaces optimization
over mixture weights. Its exact algebraic replay of all $25$ nonempty
N27/P35 mixed regions agrees with the old optimized bound. All other
old fixed-mixture bounds transfer by the pointwise domination following
Lemma~\ref{pp:mixture}; no numerical regeneration is required for this
logical transfer.
The acquisition, lower-moment, ordinary-forest and unrestricted
Fourier inputs and returned objects are retained separately in
\path{companions/finite-other-support/}. Its
\path{read_finite_other_support.py} provides the original mathematical
readers with temporary source materialization. The expensive complete
lower-moment reader is an explicit option. Partial historical count
runs retain their original domains; the analytic profile comparisons
supply the stated complementary regions.

The companion text identifies the continuous inequalities and compositions.
The frozen coefficient data, closed subdivision trees, exact arithmetic
programs and archived fresh calculations are retained in the collections
above. Their mathematical readers relate returned objects to the complete
parameter covers and consumed contracts. Fresh numerical regeneration
is a separate operation using the retained producer entrypoints.


\begin{thebibliography}{99}
\bibitem{Bentkus2004}
V. Bentkus, \emph{On Hoeffding's inequalities},
Ann. Probab. \textbf{32} (2004), 1650--1673.
\href{https://doi.org/10.1214/009117904000000360}{doi:10.1214/009117904000000360}.
\bibitem{Pinelis2006}
I. Pinelis, \emph{Binomial upper bounds on generalized moments and tail
probabilities of (super)martingales with differences bounded from above},
IMS Lecture Notes--Monograph Series \textbf{51} (2006), 33--52.
\href{https://doi.org/10.1214/074921706000000743}{doi:10.1214/074921706000000743}.
\bibitem{Schulz2016}
J. Schulz, \emph{The optimal Berry--Esseen constant in the binomial case},
Dissertation, Universit\"at Trier, 2016.
\href{https://d-nb.info/1197702695/34}{Full dissertation}.
\bibitem{Prawitz1972}
H. Prawitz, \emph{Limits for a distribution, if the characteristic
function is given in a finite domain}, Scandinavian Actuarial Journal
\textbf{1972} (1972), no.~2, 138--154.
\href{https://doi.org/10.1080/03461238.1972.10404645}{doi:10.1080/03461238.1972.10404645}.
\bibitem{Shevtsova2012}
I. Shevtsova, \emph{Moment-type estimates with asymptotically optimal
structure for the accuracy of the normal approximation},
Annales Mathematicae et Informaticae \textbf{39} (2012), 241--307.
\url{https://ami.uni-eszterhazy.hu/uploads/papers/finalpdf/AMI_39_from241to307.pdf}.
\bibitem{Shevtsova2013}
I.~G. Shevtsova, \emph{On the absolute constants in the Berry--Esseen
inequality and its structural and nonuniform improvements},
Inform. Primen. \textbf{7} (2013), no.~1, 124--125 (in Russian).
\url{https://www.mathnet.ru/eng/ia252}.
\bibitem{HeCheng2026}
H. He and G. Cheng, \emph{The Berry--Esseen Constant Conjecture is
Eventually True}, arXiv:2609.06358v3 [math.PR], 30 September 2026.
\url{https://arxiv.org/abs/2609.06358v3}.
\bibitem{RK2026}
K. Rodziewicz and B. Ko\l{}odziejek,
\emph{Moment defects and sharp iid Berry--Esseen bounds},
public working manuscript, version 3, 2026.
\href{https://hexagonmath.org/2610.00147v3}{Version record}.
\bibitem{RKFive2026}
K. Rodziewicz and B. Ko\l{}odziejek,
\emph{Five summands in the iid Berry--Esseen conjecture:
A sign-count proof and the finite-index dependency chain},
working draft, 20 October 2026.
\href{https://github.com/rodziewiczk/berry-esseen-045-certificate/blob/75a94708872e2bc053f42a279559b23ed91e4402/moment-defects/five-summands-2026-10-20/manuscript.tex}{Frozen source}.
\bibitem{RKSources2026}
K. Rodziewicz and B. Ko\l{}odziejek,
\emph{Supplementary proof and certificate sources for the iid
Berry--Esseen problem}, repository snapshot, 23 October 2026.
\href{https://github.com/rodziewiczk/berry-esseen-045-certificate/tree/75a94708872e2bc053f42a279559b23ed91e4402}{Source collection}.
\end{thebibliography}
\end{document}
